\documentclass[conference]{IEEEtran}
\IEEEoverridecommandlockouts

% \usepackage{cite}
\usepackage{amsmath,amssymb,amsfonts}
\usepackage{algorithmic}
\usepackage{graphicx}
\usepackage{textcomp}
\usepackage{xcolor}
\usepackage{url}
\usepackage{cite}
\usepackage{stfloats} 
\usepackage{booktabs}
\usepackage[bottom]{footmisc}
\usepackage{stfloats}
\usepackage[numbers, sort&compress]{natbib}
\usepackage{tabularx}

\def\BibTeX{{\rm B\kern-.05em{\sc i\kern-.025em b}\kern-.08em
    T\kern-.1667em\lower.7ex\hbox{E}\kern-.125emX}}

\makeatletter
\def\sharedaffiliation{\thanks{Contributed equally.}}
\makeatother

\begin{document}

\title{MedSwin: Multi-specialist EMR and Evidence Retrieval Architecture for Advanced Diagnosis}

\makeatletter
\newcommand{\linebreakand}{%
  \end{@IEEEauthorhalign}
  \hfill\mbox{}\par
  \mbox{}\hfill\begin{@IEEEauthorhalign}
}
\makeatother

\author{
    \IEEEauthorblockN{Dang Khoa Le\textsuperscript{*}\thanks{*Both two first authors contributed equally to this research.}}
    \IEEEauthorblockA{\textit{Department of Computing Technologies} \\
    \textit{Swinburne University of Technology} \\
    Melbourne, Australia}
    \and
    \IEEEauthorblockN{Ba Viet Anh Nguyen\textsuperscript{*}}
    \IEEEauthorblockA{\textit{Department of Computing Technologies} \\
    \textit{Swinburne University of Technology} \\
    Melbourne, Australia}
    \and
    \IEEEauthorblockN{Huu Minh Hai Truong}
    \IEEEauthorblockA{\textit{Department of Computing Technologies} \\
    \textit{Swinburne University of Technology} \\
    Melbourne, Australia}
    
    \linebreakand 

    \IEEEauthorblockN{Duy Khoa Pham}
    \IEEEauthorblockA{\textit{Department of Computing Technologies} \\
    \textit{Swinburne University of Technology} \\
    Melbourne, Australia \\
    duykhoapham@swin.edu.au}
    \and
    \IEEEauthorblockN{Bao Quoc Vo}
    \IEEEauthorblockA{\textit{Department of Computing Technologies} \\
    \textit{Swinburne University of Technology} \\
    Melbourne, Australia \\
    bvo@swin.edu.au}
    \and

}


\maketitle

\begin{abstract}
Large language models (LLMs) can generate fluent medical responses, but clinical decision support requires more than fluent text. A system must show which patient facts, guideline statements, safety constraints, and contradictory evidence support or limit an answer. Standard retrieval-augmented generation (RAG) reduces hallucination risk by adding retrieved context, but it often treats retrieval as a prompt-construction step: the top-ranked passages are passed to the generator even when they do not cover the clinical facts needed for a safe recommendation.

We present \emph{MedSwin}, an evidence-gated biomedical LLM system for patient-specific clinical decision support. The central contribution is an operational evidence-sufficiency contract. For each query, MedSwin decomposes the information need into clinical evidence facets, retrieves candidate passages from electronic medical records (EMR), clinical practice guidelines (CPGs), drug-safety sources, and biomedical literature, reranks them with a calibrated biomedical reranker, and checks whether the selected evidence bundle satisfies critical-facet coverage, provenance, contradiction, and traceability requirements before answer generation. If the evidence is incomplete or conflicting, the system retrieves more targeted evidence or produces a bounded insufficient-evidence response rather than a confident recommendation.

We evaluate MedSwin at three levels: end-to-end evidence-sufficiency audit, reranking, and standalone medical answer generation. On a TREC Clinical Decision Support system audit, MedSwin improves critical-facet recall, evidence-document recall, policy pass rate, and trace completeness under the same runtime contract. These results support the paper's main claim: clinical RAG systems should be evaluated not only by answer similarity, but by whether they retrieve sufficient evidence, know when evidence is insufficient, and expose a reviewable decision trace.
\end{abstract}

\begin{IEEEkeywords}
% • Multi-agent systems, Medical AI, Clinical decision support, Reranking, Distillation, Open-source
Multi-agent systems, Medical AI, Clinical decision support, Information retrieval, Reranking, Knowledge distillation, Large language models, Electronic medical records
\end{IEEEkeywords}


\section{Introduction}
\label{sec:intro}

Clinical decision support systems are designed to assist clinicians by connecting patient information with medical knowledge at the point of care. Evidence from controlled studies indicates that computerised decision support can improve care processes, while established reviews identify benefits in diagnostic support and guideline implementation \cite{kwan2020computerised,sutton2020overview}. Their value is increasingly important because electronic medical records (EMRs) contain extensive longitudinal information, while clinical recommendations must remain consistent with evolving clinical practice guidelines (CPGs). However, the evidence needed for an individual decision may be distributed across clinical notes and laboratory results. Conventional decision support systems also face persistent limitations, including poor workflow integration and excessive or poorly targeted alerts, which can reduce clinician acceptance and practical effectiveness \cite{sutton2020overview}.

Large language models (LLMs) offer a more flexible approach because they can interpret free text clinical information and synthesise responses to patient-specific questions. Although LLMs have achieved strong performance on medical question-answering benchmarks, such results provide limited evidence of their readiness for real world clinical decision making \cite{hager2024evaluation}. Current LLMs remain sensitive to the amount and ordering of clinical information, may fail to follow instructions, and show limited adherence to diagnostic or treatment guidelines \cite{hager2024evaluation}. Their fluent responses can also obscure whether a recommendation is supported by the relevant patient facts and current medical evidence.

Retrieval augmented generation (RAG) seeks to reduce these limitations by conditioning generation on external evidence \cite{lewis2020rag}. In medicine, RAG provides a means of connecting an LLM with current guidelines and patient specific records without relying exclusively on knowledge encoded during model training \cite{kresevic2024optimization}. However, conventional RAG pipelines typically retrieve a fixed number of highly ranked passages and place them directly into the generation context. High retrieval relevance does not ensure that the context includes every fact that could materially affect a clinical recommendation. For example, the retrieved passages may describe the appropriate treatment while omitting a patient-specific contraindication or an unresolved conflict between the EMR and the guideline.

General RAG evaluation frameworks provide limited support for examining this failure mode. RAGAS measures properties such as context relevance and answer faithfulness \cite{es2024ragas}, but these metrics do not establish whether the retrieved context contains the complete set of decision critical evidence needed for a patient-specific response. Existing medical RAG systems also tend to evaluate retrieval performance or final answer quality after a predefined context has been selected. Consequently, they provide limited control over whether generation should proceed when important evidence remains missing. They may also provide citations without exposing a structured record of which passages were considered, selected, or excluded. This restricts retrospective assessment of whether a recommendation followed from appropriate patient evidence and guideline sources. Clinical AI auditing research has emphasised the need to examine how system outputs are produced and to expose failures that aggregate performance measures may conceal \cite{liu2022medicalaudit}.

A separate challenge is adapting a compact medical language model to biomedical question answering while preserving task alignment and broader generation capabilities. SFT improves instruction following through direct optimisation on instruction-response examples \cite{wei2022finetuned}, while KD transfers teacher-generated responses and token-level predictive information to the student \cite{yang2025survey}. Both properties are valuable for biomedical question answering, where SFT aligns the model with the target task \cite{wei2022finetuned}, while KD helps retain broader teacher-derived generation behaviour \cite{yang2025survey}. To combine these complementary strengths, we first applied SFT to the KD-adapted model. However, the sequential model performed worse than the independently trained SFT and KD models, suggesting that the later SFT stage weakened capabilities acquired during distillation. This pattern is consistent with catastrophic forgetting under successive adaptation \cite{huang2024mitigating}. MedSwin therefore trains the SFT and KD branches from a shared initialization and combines their parameter updates through training-free weight-space merging \cite{yang2026model}. We evaluate Task Arithmetic, NuSLERP, DARE-Linear, and DARE-TIES to examine alternative forms of update composition and interference control \cite{yadav2023ties,yu2024language, ilharco2022editing, goddard2024arcee}.

To address these challenges, we present \emph{MedSwin}, an evidence-gated clinical RAG system implemented through centralized, orchestrator-mediated multi-agent coordination. Role-specialized components process patient-specific EMR evidence, applicable CPG evidence, and safety constraints through shared structured artefacts managed by the orchestrator. MedSwin retrieves and reranks evidence before constructing a compact bundle of selected passages and their provenance metadata. The evidence-sufficiency gate determines whether this bundle covers the required clinical facets with acceptable uncertainty and without unresolved conflicts. When the gate is not satisfied, the orchestrator performs targeted retrieval or returns a bounded insufficient-evidence response. Answer generation is permitted only after the gate accepts the evidence state, and the resulting response is accompanied by source provenance and a reviewable decision trace.

MedSwin also includes a compact medical LLM developed from a shared MedAlpaca-7B initialization. Separate SFT and KD adaptations are produced to capture complementary response behaviours. We then investigate training-free merging methods for composing these adaptations without another fine-tuning stage. This design allows the effects of SFT and KD to be evaluated independently and examines whether interference-aware merging can retain their complementary strengths at a deployable model scale.

The contributions of this work are:

\begin{enumerate}
    \item \textbf{An evidence-gated clinical decision support system.}
    We present MedSwin, a centralized, orchestrator-mediated multi-agent RAG system that applies an operational evidence-sufficiency gate before answer generation. Role-specialized components process patient-specific EMR evidence, applicable CPG evidence, and safety constraints through shared structured artefacts. The gate uses facet-level coverage, uncertainty, and contradiction checks to determine whether the system answers, performs targeted retrieval, or returns a bounded insufficient evidence response. The system also preserves an audit trace linking the final response to its supporting passages and policy decision.

    \item \textbf{A training-free method for composing medical SFT and KD behaviours.}
    We develop separate SFT-adapted and KD-adapted medical LLMs from a shared 7B initialization, then investigate whether their complementary behaviours can be combined without another gradient-based training stage. We evaluate Task Arithmetic, NuSLERP, DARE-Linear, and DARE-TIES and compare these weight-space merging methods with sequential SFT-on-KD. This study examines which merging strategy best preserves the instruction-following behaviour learnt through SFT and the teacher-derived behaviour transferred through KD.
\end{enumerate}

% Large language models (LLMs) have rapidly expanded the frontier of medical question answering and clinical text generation, but systematic reviews emphasize that reliability, hallucinations, and opaque reasoning continue to limit clinical trust and adoption \cite{meng2024llmsmedicine,kim2025medical,maynez-etal-2020-faithfulness}. In safety-critical settings, the failure mode of greatest concern is not merely an incorrect answer, but an \emph{unverifiable} answer: when the system cannot clearly expose which guideline clause, laboratory trend, or EMR event supports a recommendation, post-hoc review becomes difficult and errors can propagate. Consequently, clinical-grade systems must treat \emph{provenance, sufficiency, and auditability} as first-class design constraints rather than optional add-ons.

% Retrieval-augmented generation (RAG) partially addresses hallucinations by conditioning generation on retrieved passages from external corpora \cite{lewis2020rag}. However, most RAG pipelines remain governed by generic ranking heuristics (top-$K$ retrieval, lightweight recency/confidence filters) and evaluated with largely domain-agnostic metrics. Frameworks such as RAGAS quantify notions like context relevance and faithfulness \cite{es2024ragas}, but they do not encode clinical requirements such as guideline coverage, contraindication sensitivity, or explicit \emph{evidence sufficiency} thresholds that must be met before generating a recommendation. Recent agentic RAG assistants for eHealth environments illustrate this gap: while agentic orchestration can improve information access and report strong RAGAS-style scores, these systems often (i) rely on domain-agnostic evaluation, (ii) pass only a small set of top-ranked contexts into the prompt without calibrated inclusion policies, (iii) lack structured audit trails tied to guideline versions and EMR provenance, and (iv) depend on cloud models and crawled data that complicate privacy and governance \cite{kalathas2025agenticrag}.

% In parallel, multi-agent medical architectures have been proposed to mirror multidisciplinary clinical workflows by recruiting specialized agents (e.g., triage, differential reasoning, guideline lookup) and coordinating them through an orchestrator \cite{sulis2023agents,kim2024mdagents}. While this decomposition improves modularity and can reduce single-model brittleness, most existing designs still treat retrieval as a best-effort prelude rather than a controlled, auditable decision stage. In clinical RAG, small ranking errors can omit critical contraindications or include irrelevant background, and uncalibrated relevance scores make threshold-based policies difficult to justify. Moreover, continual updates via sequential fine-tuning can overwrite previously acquired behaviors, producing instability and catastrophic forgetting—an undesirable property for systems intended to operate under fixed institutional policies \cite{luo2023catastrophic}.

% \textbf{This work.} We present \emph{MedSwin}, a biomedical LLM system for patient-specific clinical question answering. The core idea is simple: before the model answers, it must first retrieve enough relevant evidence from the patient's EMR and the relevant clinical guidelines. If key evidence is missing, unclear, or contradictory, the system should not produce a confident recommendation. It should either retrieve more evidence, narrow the answer, or state that the available evidence is insufficient.

% MedSwin therefore treats evidence selection as a required step before generation, rather than as an optional explanation after generation. The system first retrieves candidate EMR and CPG passages, reranks them with a biomedical reranker, checks whether the selected evidence covers the critical clinical facets of the question, and only then asks the medical LLM to generate a grounded answer. In parallel, we train a compact medical LLM through supervised fine-tuning, knowledge distillation, and training-free model merging so that the final system remains deployable while preserving medical reasoning and response discipline.

% Concretely, we make the following contributions:
% \begin{enumerate}
%     \item \textbf{An evidence-gated clinical decision-support system.}
%     We introduce a patient-specific clinical RAG architecture that treats evidence sufficiency as a precondition for answer generation. MedSwin combines biomedical retrieval with calibrated reranking, then evaluates facet coverage and unresolved contradiction before permitting synthesis. A bounded multi-agent coordination layer performs targeted evidence exploration when a critical facet is missing. The system produces a traceable evidence bundle that links the final response to its supporting EMR and CPG passages.

%     \item \textbf{A training-free method for composing medical SFT and KD behaviours.}
%     We develop separate SFT-adapted and KD-adapted medical LLMs from a shared 7B initialization, then study weight-space merging as an alternative to sequential fine-tuning. The evaluated methods include TIES and DARE-based composition, which are designed to reduce interference between task-vector updates. This study examines whether instruction-following discipline and teacher-derived semantic behaviour can be combined without additional gradient-based training.
% \end{enumerate}

\section{Related Work}
\label{sec:related_work}

% MedSwin builds on medical LLMs, biomedical retrieval, RAG evaluation, and multi-agent clinical decision support. The closest gap is not the absence of RAG itself, but the absence of a clinical RAG control layer that decides whether retrieved evidence is sufficient before generation.

% \paragraph{Medical LLMs}
% Medical LLM research has improved biomedical question answering and clinical text generation through domain pretraining, instruction tuning, and larger medical foundation models \cite{singhal2023medpalm2,chen2023meditron,labrak2024biomistral,han2023medalpaca,li2023chatdoctor,google2025medgemma}. These models are valuable generators, but answer quality alone does not establish whether a patient-specific response is supported by the right EMR and guideline evidence. MedSwin therefore treats the medical LLM as the final synthesiser, not as the sole source of clinical reliability.

% \paragraph{Medical RAG and RAG evaluation}
% RAG grounds generation in retrieved documents and reduces reliance on parametric memory \cite{lewis2020rag}. RAG evaluation frameworks such as RAGAS measure context relevance, answer faithfulness, and related retrieval-generation signals \cite{es2024ragas}. These metrics are important, but clinical decision support also requires patient applicability, contraindication coverage, guideline-version provenance, and explicit handling of missing or conflicting evidence. MedSwin addresses this gap by making evidence sufficiency a pre-generation acceptance condition.

% \paragraph{Biomedical retrieval and reranking}
% Biomedical retrieval models such as MedCPT show that domain-specific retrieval improves biomedical search \cite{medcpt}. Two-stage retrieval pipelines with neural reranking are effective for high-precision passage selection \cite{nogueira2019passage_bert}. Open LLM-based rerankers such as BGE-reranker-v2-gemma make long-context reranking more accessible \cite{bgev2gemma}. Open, long-context biomedical rerankers with calibrated, policy-usable scores suitable for auditable EMR/CPG evidence selection remain limited in the literature. MedSwin extends this pattern by calibrating biomedical reranker scores and using them inside explicit evidence-selection policies.

% \paragraph{Multi-agent clinical systems}
% Multi-agent healthcare systems can decompose complex workflows into specialised roles, and recent LLM-agent work has explored collaborative medical reasoning and agentic RAG \cite{sulis2023agents,kim2024mdagents,choi2024malade,kalathas2025agenticrag,wang2025llmagents}. However, agentic orchestration can still fail if the agents produce ungrounded summaries or if retrieval provenance is treated as secondary. MedSwin uses agents to search for and verify missing clinical facets, while the final answer remains controlled by the evidence-sufficiency contract.

% \paragraph{Positioning}
% MedSwin differs from prior medical RAG and agentic medical assistants in one central respect: the system does not treat retrieved context as automatically usable. It first checks whether the evidence bundle is sufficient for the clinical question. The contribution is therefore not simply a new RAG pipeline, a new reranker, or a new medical LLM. It is the integration of calibrated retrieval, facet-level sufficiency checking, contradiction handling, and audit-trace generation into one acceptance-controlled clinical RAG system.


% -
% Figure
% -
% \begin{figure}[h!]
%     \centering
%     \includegraphics[width=1\linewidth]{intro.png}
%     \caption{MedSwin evidence-gated clinical RAG system. Retrieval and reranking produce evidence bundle; sufficiency gate decides whether the system answers, retrieves more evidence, or returns insufficient-evidence response.}
%     \label{fig:intro-diagram}
% \end{figure}


\section{MedSwin Models}
% The MedSwin language model is developed from MedAlpaca-7B through two independent adaptation pathways. The first uses supervised fine-tuning (SFT) to strengthen biomedical instruction following. The second uses knowledge distillation (KD) to transfer information from MedGemma-27B-it. Both checkpoints retain the same base initialization, which enables their parameter updates to be composed through training-free model merging. This design allows SFT, KD, sequential adaptation, and weight-space composition to be evaluated under a controlled model scale.

\subsection{Supervised Fine-Tuning}
\label{sec:medswin-sft}

% We first adapt MedAlpaca-7B with supervised fine-tuning (SFT) on a mixed biomedical instruction corpus covering patient-facing questions, clinician-style explanations, and evidence-oriented biomedical QA. The corpus combines HealthCareMagic, iCliniq, and PubMedQA, with medically constrained augmentation used to improve linguistic coverage rather than to introduce new clinical facts.

% \begin{figure}[h!]
%     \centering
%     \includegraphics[width=0.9\linewidth]{aug.png}
%     \caption{Biomedical QA data augmentation pipeline for SFT.}
%     \label{fig:augmentation-pipeline}
% \end{figure}

% Each example is serialized as an instruction-input-output triple $(x_i,y_i)$, and the student is trained with token-level cross-entropy on the target completion:
% \begin{equation}
% \mathcal{L}_{\mathrm{SFT}}
% =
% -\sum_{t=1}^{|y_i|}
% \log p_{\theta}(y_{i,t}\mid x_i,y_{i,<t}).
% \end{equation}
% To reduce source-style overfitting, we stratify the data mixture across corpora and augmentation types, remove near-duplicates, and reject augmented samples that alter medical meaning, including dosage, contraindication, temporal, or uncertainty statements. Full data composition, augmentation prompts, and quality-control rules are provided in Appendix~\ref{app:sft-data-augmentation}.

% We adapt MedAlpaca-7B through supervised fine-tuning using QLoRA \cite{dettmers2023qlora}. The training corpus combines HealthCareMagic, iCliniq, and PubMedQA to provide patient-facing medical questions and evidence-oriented biomedical question answering \cite{han2023medalpaca,li2023chatdoctor,pubmedqa}. As shown in Fig.~\ref{fig:augmentation-pipeline}, the source examples pass through a controlled construction pipeline that expands linguistic variation and removes samples that alter the medical meaning of the original content. Detailed data composition and validation procedures are provided in Appendix~\ref{app:sft-data-augmentation}.

% \begin{figure}[t]
%     \centering
%     \includegraphics[width=\linewidth]{aug.png}
%     \caption{Construction of the biomedical instruction corpus used for SFT.
%     Source datasets are normalised and validated before augmentation through
%     back-translation or synthetic generation. PHI scrubbing, deduplication,
%     and quality assurance are applied before accepted examples are written to
%     the final corpus.}
%     \label{fig:augmentation-pipeline}
% \end{figure}

% Each example is represented as an instruction-response pair $(x_i,y_i)$, where $x_i$ contains the instruction and any associated context, and $y_i$ is the target response. The model is trained by minimising the token-level cross-entropy loss

% \begin{equation}
% \mathcal{L}_{\mathrm{SFT}}
% -\frac{1}{N_{\mathrm{ans}}}
% \sum_{i=1}^{N}
% \sum_{t=1}^{|y_i|}
% \log p_{\theta}
% \left(
% y_{i,t} \mid x_i, y_{i,<t}
% \right),
% \label{eq:sft-loss}
% \end{equation}

% where $y_{i,t}$ denotes the target token at position $t$, $y_{i,<t}$ denotes the preceding target tokens, and $N_{\mathrm{ans}} = \sum_{i=1}^{N} |y_i|$ represents the total number of target answer tokens across all $N$ examples.

% To reduce the computational cost of adapting the 7B model, we employ QLoRA, which backpropagates gradients through a frozen 4-bit quantised backbone into trainable low-rank adapters \cite{dettmers2023qlora}. QLoRA was shown to retain full 16-bit fine-tuning task performance while substantially reducing memory requirements \cite{dettmers2023qlora}, making it suitable for adapting MedAlpaca-7B under limited GPU resources. Gradients are propagated through the quantised backbone into trainable low-rank adapters attached to selected transformer projections. This QLoRA configuration avoids storing full-precision optimizer states for all model parameters and enables biomedical instruction tuning under limited GPU memory, while preserving the pretrained backbone for the subsequent controlled comparison with the KD branch \cite{dettmers2023qlora}.

% The resulting adapter is merged with the shared MedAlpaca-7B initialization to produce the SFT checkpoint $\theta_{\mathrm{SFT}}$. This checkpoint is trained independently from $\theta_{\mathrm{KD}}$, allowing the parameter updates induced by SFT and KD to be compared and composed through the training-free merging strategies described in Section~\ref{sec:merge}.

We adapt MedAlpaca-7B through supervised fine-tuning using QLoRA \cite{dettmers2023qlora} to reduce the computational cost of adapting the 7B model. QLoRA backpropagates gradients through a frozen 4-bit quantised backbone into trainable low-rank adapters and was shown to retain full 16-bit fine-tuning task performance while substantially reducing memory requirements \cite{dettmers2023qlora}, making it suitable for adapting MedAlpaca-7B under limited GPU resources. Gradients are propagated through the quantised backbone into trainable low-rank adapters attached to selected transformer projections. This configuration avoids storing full-precision optimiser states for all model parameters and enables biomedical instruction tuning under limited GPU memory, while preserving the pretrained backbone for the subsequent controlled comparison with the KD branch \cite{dettmers2023qlora}.

The training corpus combines HealthCareMagic, iCliniq, and PubMedQA to provide patient-facing medical questions and evidence-oriented biomedical question answering \cite{han2023medalpaca,li2023chatdoctor,pubmedqa}. As shown in Fig.~\ref{fig:augmentation-pipeline}, the source examples pass through a controlled construction pipeline that expands linguistic variation and removes samples that alter the medical meaning of the original content. Detailed data composition and validation procedures are provided in Appendix~\ref{app:sft-data-augmentation}.

\begin{figure}[t]
    \centering
    \includegraphics[width=\linewidth]{aug.png}
    \caption{Construction of the biomedical instruction corpus used for SFT.
    Source datasets are normalised and validated before augmentation through
    back-translation or synthetic generation. PHI scrubbing, deduplication,
    and quality assurance are applied before accepted examples are written to
    the final corpus.}
    \label{fig:augmentation-pipeline}
\end{figure}

Each example is represented as an instruction-response pair $(x_i,y_i)$, where $x_i$ contains the instruction and any associated context, and $y_i$ is the target response. The model is trained by minimising the token-level cross-entropy loss

\begin{equation}
\mathcal{L}_{\mathrm{SFT}}
=
-\frac{1}{N_{\mathrm{ans}}}
\sum_{i=1}^{N}
\sum_{t=1}^{|y_i|}
\log p_{\theta}
\left(
y_{i,t} \mid x_i, y_{i,<t}
\right),
\label{eq:sft-loss}
\end{equation}

where $y_{i,t}$ denotes the target token at position $t$, $y_{i,<t}$ denotes the preceding target tokens, and $N_{\mathrm{ans}}=\sum_{i=1}^{N}|y_i|$ represents the total number of target answer tokens across all $N$ examples. 

The resulting adapter is merged with the shared MedAlpaca-7B initialisation to produce the SFT checkpoint $\theta_{\mathrm{SFT}}$. This checkpoint is trained independently from $\theta_{\mathrm{KD}}$, allowing the parameter updates induced by SFT and KD to be compared and composed through the training-free merging strategies described in Section~\ref{sec:merge}.

\subsection{Knowledge Distillation}
\label{sec:medswin-kd}

% While SFT equips the student with instruction-following behaviour and biomedical task competence, we further transfer domain-specific reasoning from MedGemma-27B-it into MedAlpaca-7B through knowledge distillation. Our objective is to improve answer quality and preserve uncertainty-aware behaviours that are important in medical decision support.

% We employs a hybrid distillation strategy that combines teacher-generated responses (hard-label distillation) with teacher probability distributions (soft-label distillation). Hard labels expose the student to a larger and more diverse set of biomedical question-answering patterns, while soft labels transfer richer information about alternative plausible responses and uncertainty structure. The resulting student model remains deployable at 7B scale while inheriting clinical reasoning behaviours from a substantially larger medical instructor. Detailed data-flow, optimisation objectives, training settings, and implementation details are provided in Appendix~\ref{app:kd-dataflow} and Appendix~\ref{app:kd-practicals}.

% \paragraph{Positioning and gap} Prior open medical instruction-tuned LLMs largely emphasize supervised instruction finetuning on curated or mined medical dialogues and QA corpora \cite{han2023medalpaca,li2023chatdoctor}. In contrast, systematic transfer of uncertainty-aware behaviours from a stronger medical instructor via token-distribution matching remains less explored in open medical pipelines. MedSwin therefore investigates whether combining hard-label expansion, soft-label distillation, and parameter-efficient adaptation can produce a compact medical specialist that retains both medical knowledge and calibrated response behaviour.

We construct a second MedAlpaca-7B adaptation through supervised knowledge distillation from MedGemma-27B-it \cite{han2023medalpaca,google2025medgemma}. This branch is trained independently from the SFT checkpoint and uses the same MedAlpaca-7B initialisation. The purpose is to expose the student to teacher generated biomedical responses and token-level predictive information that is absent from one hot instruction targets.

MedSwin combines hard label and soft label distillation \cite{hinton2015distilling}. For hard label distillation, MedGemma-27B-it generates responses to unlabelled PubMedQA questions and synthesised biomedical prompts. These prompt-response pairs are treated as fixed instruction targets and added to the biomedical training corpus. For soft label distillation, the student matches the teacher's next-token distribution under the same teacher-forced response prefix. These soft targets retain information about alternative token choices that is lost when training only on the generated response.

To reduce the storage cost of offline distillation, MedSwin retains only the teacher's top-$k$ token log-probabilities at each aligned answer step. Prior LLM distillation work has shown that a concentrated subset of teacher (top-k) probabilities can support efficient knowledge transfer \cite{shum2024first}. Let $K_{b,t}$ denote the retained token set for batch item $b$ at answer step $t$. The stored teacher log-probabilities $\ell_{T,b,t,j}$ are renormalised over $K_{b,t}$ as

\begin{equation}
\widetilde{p}_{T}(j\mid b,t)
=
\frac{\exp\!\left(\ell_{T,b,t,j}\right)}
{\sum\limits_{m\in K_{b,t}}
 \exp\!\left(\ell_{T,b,t,m}\right)},
\qquad j\in K_{b,t}.
\label{eq:kd-teacher-topk}
\end{equation}

MedSwin reconstructs the teacher distribution from the stored top-$k$ log-probabilities and applies temperature scaling to the student logits over the same retained token set:
\begin{equation}
p_{S}^{(\tau)}(j\mid b,t)
=
\frac{\exp\!\left(z_{S,b,t,j}/\tau\right)}
{\sum\limits_{m\in K_{b,t}}
 \exp\!\left(z_{S,b,t,m}/\tau\right)},
\qquad j\in K_{b,t},
\label{eq:kd-student-topk}
\end{equation}

where $\tau$ is the KD temperature.

The hard label term is the cross-entropy loss computed only over answer tokens:

\begin{equation}
\mathcal{L}_{\mathrm{CE}}
=
-\frac{1}{N_{\mathrm{ans}}}
\sum_{b,t}
m_{b,t}
\log p_{S}
\left(
y_{b,t}\mid x_b,y_{b,<t}
\right),
\label{eq:kd-ce}
\end{equation}

where $m_{b,t}\in\{0,1\}$ masks prompt and padding positions, and $N_{\mathrm{ans}}=\sum_{b,t}m_{b,t}$. The soft label term applies forward KL divergence at positions with valid aligned teacher targets:

\begin{equation}
\mathcal{L}_{\mathrm{soft}}
=
\frac{1}{|\Omega|}
\sum_{(b,t)\in\Omega}
D_{\mathrm{KL}}
\left(
\widetilde{p}_{T}(\cdot\mid b,t)
\,\middle\|\,
p_{S}^{(\tau)}(\cdot\mid b,t)
\right),
\label{eq:kd-soft}
\end{equation}

where $\Omega$ is the set of valid aligned KD steps. The total training loss is

\begin{equation}
\mathcal{L}_{\mathrm{KD}}
=
\mathcal{L}_{\mathrm{CE}}
+
\lambda\mathcal{L}_{\mathrm{soft}},
\label{eq:kd-total}
\end{equation}

where $\lambda$ controls the contribution of the soft label term.

Student adaptation is performed with QLoRA \cite{dettmers2023qlora}. The MedAlpaca-7B backbone is loaded with 4-bit NF4 quantisation and prepared for $k$-bit training, while only the low-rank LoRA adapter parameters are updated. Gradient checkpointing is enabled to reduce memory usage, and early stopping is based on validation loss over held-out biomedical QA splits. The resulting KD adapter checkpoint is saved separately from the SFT adapter checkpoint for the training-free composition described in Section~\ref{sec:merge}. Full data-alignment and optimisation settings are provided in Appendices~\ref{app:kd-dataflow} and \ref{app:kd-practicals}.

\subsection{Training-Free Model Merging}
\label{sec:merge}

% Supervised fine-tuning (SFT) and knowledge distillation (KD) encode partially complementary inductive biases into the same base model, yet their parameter updates can interfere. SFT optimizes next-token likelihood on instruction-response pairs, sharpening format adherence, response discipline, and stylistic alignment. KD, by matching the teacher distribution, can transfer calibrated uncertainty and semantic behaviors that are not fully expressed by hard labels alone \cite{hinton2015distilling}. In safety-critical biomedical QA, this tension is practical: SFT often improves conservatism and output structure, whereas KD can preserve teacher-like semantic fidelity and hedging patterns that reduce overconfident error modes.

% \paragraph{Goal and contribution}
% Our objective is a compact medical assistant that simultaneously (i) preserves teacher-like semantic fidelity and uncertainty behaviors transferred by KD, and (ii) retains the conservative formatting and response discipline induced by SFT. We contribute a \emph{training-free} composition study for biomedical QA that: (a) formalizes SFT and KD adaptations as task vectors from a shared initialization; (b) evaluates merge operators with explicit interference-control mechanisms (DARE, TIES, and their composition) \cite{yu2023supermario,yadav2023ties}; and (c) empirically contrasts weight-space merging against sequential SFT-on-KD, which is known to trigger catastrophic forgetting under continual fine-tuning/instruction tuning \cite{luo2023catastrophic}. Our perspective is consistent with the emerging view that fine-tuned solutions from a common base can remain within broad, low-error basins where weight-space composition (when carefully controlled) can be beneficial \cite{wortsman2022soups,yang2026model}. We implement merging using an established toolkit to ensure reproducibility and alignment with common open-source workflows \cite{goddard2024a}.

% \paragraph{Notation (task vectors)}
% Let $\theta_0$ denote the shared initialization (MedAlpaca-7B), $\theta_{\text{SFT}}$ the SFT-adapted checkpoint, and $\theta_{\text{KD}}$ the KD-adapted checkpoint. We define task vectors (parameter deltas)
% \begin{equation}
% \Delta_{\text{SFT}} = \theta_{\text{SFT}} - \theta_0, \qquad
% \Delta_{\text{KD}} = \theta_{\text{KD}} - \theta_0.
% \end{equation}
% Training-free merging constructs $\theta_{\text{merge}}$ directly in weight space (no additional SGD), which avoids gradient overwriting and the associated forgetting risk \cite{luo2023catastrophic}. Conceptually, merging aims to preserve behaviors that are represented by non-conflicting (or weakly conflicting) update directions while suppressing destructive parameter interference \cite{yadav2023ties,yu2023supermario}.

% \subsubsection*{Combination strategies}

% \paragraph{Sequential SFT-on-KD}
% A naive strategy is to continue instruction fine-tuning from $\theta_{\text{KD}}$ using the SFT objective, producing $\theta_{\text{KD}\rightarrow\text{SFT}}$ (1). While this can improve stylistic alignment, empirical evidence shows that continual fine-tuning often overwrites previously acquired capabilities due to gradient interference and representational drift, especially for 1B-7B LLMs \cite{luo2023catastrophic}. In our setting, this risk is amplified because SFT pushes toward narrower stylistic/format manifolds that may not preserve distributional behaviors transferred via KD.

% \paragraph{Weight-space merging (training-free)}
% We instead merge independently trained SFT and KD checkpoints without additional optimization (2). This is motivated by: (i) observations that averaging/interpolating fine-tuned checkpoints from a shared base can improve robustness without extra inference cost \cite{wortsman2022soups}; and (ii) task-vector arithmetic enabling explicit control over how much of each behavioral update is injected into the merged model \cite{ilharco2022editing}. We focus on interference-aware merges because medical SFT and KD deltas are often broad and partially misaligned, making naive averaging prone to destructive cancellation.

% \paragraph{Merge operators}
% We evaluate Task Arithmetic, NuSLERP, DARE-Linear, TIES, and DARE-TIES as training-free operators for composing SFT and KD task vectors. These operators test increasingly explicit forms of interference control: magnitude scaling, direction interpolation, stochastic sparsification, sign-consensus filtering, and their composition. Full operator definitions and equations are provided in Appendix~\ref{app:merge-operators}.

% \section{MedSwin Architecture}

% \subsection{System Overview}
% \label{subsec:techniques}

% MedSwin answers a clinical question in four steps. First, it retrieves candidate passages from the patient's EMR, relevant clinical practice guidelines (CPG), drug-safety resources, and biomedical literature. Second, a biomedical reranker scores each passage for relevance to the question. Third, the system checks whether the selected passages cover the critical clinical facets needed for a safe answer. Finally, the medical LLM generates an answer using only the accepted evidence bundle and its provenance metadata.

% The system produces three outputs for each query: (i) a clinical answer, (ii) an evidence bundle containing the EMR and CPG passages used to support the answer, and (iii) an audit trace showing which passages were retrieved, reranked, selected, rejected, or flagged as insufficient. This makes evidence sufficiency a precondition for answer generation rather than a post-hoc explanation.

SFT and KD are combined because they contribute different capabilities needed for biomedical question answering. The SFT branch adapts MedAlpaca-7B to the target instruction corpus, strengthening instruction following and response alignment. The KD branch transfers teacher-generated responses and token-level predictive information from MedGemma-27B-it, helping the student retain broader teacher-derived generation behaviour. A compact medical model should preserve both properties: task alignment from SFT and broader generation capability from KD. Applying SFT directly to the KD-adapted model provides a straightforward way to combine them, but the later optimisation stage can overwrite part of the distilled adaptation. MedSwin therefore trains the SFT and KD branches independently from the same MedAlpaca-7B initialization and combines their parameter updates through training-free weight-space merging. This design allows the two adaptations to be composed without exposing either checkpoint to another gradient-based training stage.

\paragraph{Checkpoint construction}
Let $\theta_0$ denote the shared MedAlpaca-7B initialization. The SFT and KD branches are adapted independently using parameter-efficient fine-tuning. Before model merging, each trained adapter is folded into a separate copy of $\theta_0$, producing the full SFT and KD checkpoints $\theta_{\mathrm{SFT}}$ and $\theta_{\mathrm{KD}}$. This materialisation step places both adaptations in the same parameter space and preserves a common reference model for task-vector construction.

The parameter updates induced by the two branches are defined as

\begin{equation}
\Delta_{\mathrm{SFT}}
=
\theta_{\mathrm{SFT}}-\theta_0,
\qquad
\Delta_{\mathrm{KD}}
=
\theta_{\mathrm{KD}}-\theta_0.
\label{eq:merge-task-vectors}
\end{equation}

A training-free merged checkpoint is then constructed as

\begin{equation}
\theta_{\mathrm{merge}}
=
\theta_0+
\mathcal{M}
\left(
\Delta_{\mathrm{SFT}},
\Delta_{\mathrm{KD}};\boldsymbol{\lambda}
\right),
\label{eq:general-model-merge}
\end{equation}

where $\mathcal{M}$ denotes the merging operator and $\boldsymbol{\lambda}$ contains its scaling or interpolation parameters. This formulation combines independently learned adaptations without an additional gradient-based training stage \cite{ilharco2022editing}. The merged model retains the architecture and parameter count of the 7B base model, so merging introduces no ensemble-related inference cost.

\paragraph{Sequential adaptation baseline}

We use sequential SFT-on-KD as the gradient-based composition baseline. Starting from $\theta_{\mathrm{KD}}$, the model is further adapted using the SFT corpus and the loss in Eq.~\eqref{eq:sft-loss}, producing

\begin{equation}
\theta_{\mathrm{KD}\rightarrow\mathrm{SFT}}
=
\operatorname{SFT}
\left(
\theta_{\mathrm{KD}}
\right).
\label{eq:sequential-kd-sft}
\end{equation}

This baseline tests whether continued optimisation is more effective than combining the independently learned SFT and KD updates. It also provides a controlled comparison between sequential parameter modification and training-free weight-space composition.

\paragraph{Merge operators}

We evaluate four operators that differ in how they combine update magnitude and direction. Task Arithmetic provides the direct linear baseline by adding scaled SFT and KD task vectors \cite{ilharco2022editing}. NuSLERP performs normalised spherical interpolation between the two update directions, reducing sensitivity to differences in their scale; we use the implementation provided by MergeKit \cite{goddard2024arcee}. DARE-Linear first applies stochastic drop-and-rescale sparsification to each task vector and then performs linear composition \cite{yu2024language}. DARE-TIES applies the same sparsification before sign-aware TIES merging, which trims low-magnitude updates and resolves coordinate-level sign disagreement among the retained parameters \cite{yu2024language,yadav2023ties}.

These operators provide a controlled progression from direct linear composition to geometry-aware interpolation and explicit interference management. Task Arithmetic tests whether the two adaptations can be combined without specialised conflict handling. NuSLERP examines whether directional interpolation is sufficient, while the two DARE variants assess the effect of sparsification with and without sign-consensus resolution. Full operator definitions are provided in Appendix~\ref{app:merge-operators}.

All merge coefficients, interpolation weights, and sparsification rates are selected using validation data and fixed before test-set evaluation. Every operator is applied to the same SFT and KD checkpoints, ensuring that observed differences can be attributed to the composition strategy instead of differences in training data or model initialization.

\section{MedSwin Evidence-Gated Clinical RAG System}
\label{sec:medswin-system}

\subsection{System Overview}
\label{sec:system-overview}
MedSwin is an evidence-gated clinical RAG system implemented through a centralized, orchestrator-mediated multi-agent architecture. It receives a clinical query and a patient identifier used to retrieve the corresponding EMR context. A single orchestrator invokes role-specialized agents in a fixed workflow and transfers their outputs through structured artefacts. The agents do not communicate directly or negotiate through a peer-to-peer protocol.

Figure~\ref{fig:medswin-system} presents the end-to-end runtime. The supervisor role normalises the query and produces a structured retrieval specification. Hybrid retrieval obtains candidate passages from patient-specific EMRs and external clinical sources, including CPGs. The reranker then orders these passages before the orchestrator constructs a compact evidence bundle. The resulting evidence state is processed through role-specific EMR, guideline, and safety functions before the evidence-sufficiency gate determines whether generation is permitted.

The gate produces one of three outcomes: \textsc{Answer}, \textsc{Retrieve-More}, or \textsc{Insufficient-Evidence}. Under \textsc{Answer}, the accepted evidence bundle and its structured summaries are passed to the MedSwin language model for final composition. Under \textsc{Retrieve-More}, the orchestrator issues a targeted retrieval request for a missing or uncertain facet. When the gate cannot be satisfied within the available sources or retrieval-loop budget, the system returns \textsc{Insufficient-Evidence} and limits the response accordingly.

Final answer generation occurs only after the evidence decision, while MedSwin language model step normalizes the query for retrieval. It generates the final response from the accepted evidence bundle and its provenance metadata. For each query, the system returns either a grounded clinical answer or a bounded insufficient-evidence response. It also preserves the supporting evidence and a structured audit trace recording retrieval, reranking, evidence selection, and the final sufficiency decision.

\begin{figure}[t]
    \centering
    \includegraphics[width=\linewidth]{intro.png}
    \caption{Overview of the MedSwin evidence-gated clinical RAG system. Patient-specific EMR and CPG evidence is retrieved and reranked before the evidence-sufficiency gate determines whether the system generates an answer, retrieves additional evidence, or returns an insufficient-evidence response. Accepted evidence and provenance metadata are passed to the MedSwin language model to produce a grounded answer with a reviewable audit trace.}
    \label{fig:medswin-system}
\end{figure}

\subsection{Clinical Evidence Facets}
\label{sec:clinical-facets}

% Query level relevance alone does not establish that the retrieved evidence contains every clinical fact needed to support a patient-specific response. A set of highly ranked passages may discuss the relevant treatment while omitting a patient constraint, an applicable guideline condition, or evidence needed to interpret the patient's current state. Hence, MedSwin represents the evidence requirements of each clinical query as a set of distinct clinical evidence facets.

% For a clinical query \(q\), the required facet set is denoted by
% \[
% \mathcal{F}(q)=\{f_1,\ldots,f_m\},
% \]

% where \(q\) is the clinical query and \(f_j\) denotes the \(j\)-th decision-relevant information need. For example, a medication safety query may require evidence about the patient's current medication and recent laboratory status. It may also require the applicable guideline condition or a relevant safety restriction.

% MedSwin further identifies a decision-critical subset:
% \[
% \mathcal{F}_{\mathrm{critical}}(q)\subseteq\mathcal{F}(q).
% \]

% A critical facet contains information whose absence could affect the validity or safety of the system response. In the current system design, such facets include  evidence of treatment recommendation, contraindications, adverse-risk evidence, and patient applicability.

% For each candidate passage \(d\) and facet \(f_j\in\mathcal{F}(q)\), MedSwin uses a facet-alignment score \(a_{dj}\in[0,1]\), defined as
% \[
% a_{dj}
% =
% P\!\left(
% f_j \text{ is adequately addressed by passage } d
% \mid q
% \right),
% \]

% where \(d\) denotes a candidate evidence passage. A higher \(a_{dj}\) indicates that the passage provides stronger evidence for the corresponding facet in the context of the query. This score contributes to the estimated coverage of each facet during evidence-bundle construction.

% The facet representation connects evidence retrieval with the subsequent sufficiency decision. During bundle construction, passages that cover previously unsupported facets receive greater utility. The resulting facet coverage state is then examined by the evidence-sufficiency gate, with stricter requirements applied to \(\mathcal{F}_{\mathrm{critical}}(q)\). When a critical facet remains missing or uncertain, the system may retrieve additional evidence or return a bounded insufficient-evidence response.

Query level relevance does not establish that the retrieved evidence contains the clinical information required for a patient-specific response. Several passages may be strongly related to the overall topic while omitting a contraindication or a condition of patient applicability that could change the system decision. Hence, MedSwin decomposes each query into clinical evidence facets before evidence selection.

For a clinical query \(q\), the resulting facet set is denoted by
\[
\mathcal{F}(q)=\{f_1,\ldots,f_m\},
\]
where \(m\) is the number of identified facets and \(f_j\) denotes the \(j\)-th evidence requirement, for \(j\in\{1,\ldots,m\}\). The orchestrator first normalises the query using a structured LLM prompt that extracts relevant medical terms and retrieval hints, then converts the required evidence needs into facet records. Each facet \(f_j\) includes a coverage threshold \(\theta_j\), which defines the minimum evidence coverage required for that facet, and an importance weight \(w_j\), which controls its contribution during evidence selection. The facet also specifies its expected evidence source and retrieval keywords. If query normalisation does not produce usable facets, MedSwin applies a predefined fallback set covering guideline concordance, safety constraints, patient applicability, and evidence quality.

Let \(r_j\in\{0,1\}\) indicate whether facet \(f_j\) is required for the current query. The decision-critical facet set is therefore defined as
\[
\mathcal{F}_{\mathrm{critical}}(q)
=
\{f_j\in\mathcal{F}(q)\mid r_j=1\}.
\]
Evidence for these facets must satisfy their corresponding coverage thresholds before the system can permit answer generation.

For each candidate passage \(d\), MedSwin computes a bounded facet-alignment score 
\(a_{dj}\in[0,1]\):
\begin{equation}
\label{eq:facet-alignment}
\small
\begin{aligned}
a_{dj}
=
\operatorname{clip}_{[0,1]}
\Big[
&s_{\mathrm{src}}(d,f_j)
+s_{\mathrm{key}}(d,f_j) \\
&+s_{\mathrm{sec}}(d)
+s_{\mathrm{safety}}(d,f_j)
\Big].
\end{aligned}
\end{equation}
where \(s_{\mathrm{src}}\) measures compatibility between the passage source and the facet source policy, \(s_{\mathrm{key}}\) measures support from matched facet keywords, and \(s_{\mathrm{sec}}\) represents the informativeness of the passage section. The term \(s_{\mathrm{safety}}\) contributes additional alignment when a passage contains safety evidence relevant to a contraindication facet. When none of these signals provides support, the alignment score is set to zero.

During evidence selection, passages receive greater utility when they address highly weighted facets that are not yet covered by the current evidence bundle. The resulting passage-facet alignment scores are then used to estimate facet coverage and support the evidence-sufficiency decision.

\subsection{Hybrid Retrieval and Reranking}
\label{sec:retrieval-reranking}
Retrieving more passages does not necessarily provide a stronger basis for answer generation. A larger candidate set may improve evidence recall, but it can also introduce irrelevant passages that reduce the generator's ability to use the important information \cite{yu2024rankrag,liu2024lost}. Therefore, MedSwin separates broad candidate retrieval from precise reranking. The first stage retrieves potentially useful evidence, while the second stage prioritises passages that are most relevant to the clinical query.

MedSwin searches patient-specific EMRs and external clinical sources, including CPGs. Documents are divided into section-aware passages before indexing. Each passage retains its document identifier and source, together with section and temporal metadata, to support patient scoping and subsequent provenance tracing.

\paragraph{Hybrid candidate retrieval}
MedSwin combines dense and lexical retrieval. Dense retrieval represents the clinical query and indexed passages as embeddings, then searches for passages that are semantically close to the query. The passage embeddings are stored in an HNSW approximate nearest-neighbour index configured with cosine distance. For a query \(q\) and indexed passage \(d\), the returned distance is converted to a dense similarity score as
\begin{equation}
\label{eq:dense-retrieval}
s_{\mathrm{dense}}(q,d)
=
1-\operatorname{dist}_{\mathrm{cos}}(q,d).
\end{equation}

HNSW returns the dense candidate set \(\mathcal{C}_{\mathrm{dense}}(q)\). In parallel, BM25 retrieves the lexical candidate set \(\mathcal{C}_{\mathrm{lex}}(q)\), preserving exact matches for clinical terminology and abbreviations. The two sets are combined to form the reranking pool:
\begin{equation}
\label{eq:hybrid-candidates}
\small
\mathcal{C}(q)
=
\operatorname{TopK}_{K'}
\left(
\mathcal{C}_{\mathrm{dense}}(q)
\cup
\mathcal{C}_{\mathrm{lex}}(q)
\right),
\end{equation}
where \(K'\) is the maximum number of candidates retained for reranking. Metadata filters may restrict candidates by patient identifier, source, section, or time when required by the query.

\paragraph{Reranking}
Each query-passage pair \((q,d)\in\mathcal{C}(q)\) is evaluated independently by the biomedical reranker. For each candidate passage, the reranker returns a relevance estimate
\[
\widehat{p}_{\theta}(q,d)\in[0,1],
\]
together with optional scoring metadata. MedSwin stores this value as the passage-level relevance score and preserves the reranked order of the candidate pool.

The reranker is used to refine the broad hybrid-retrieval output before evidence fusion. Its relevance estimate contributes to the subsequent fusion score and to the confidence assigned to passage-facet evidence. Final bundle selection is based on the combined evidence signals described in the next subsection, rather than on the reranker score alone. Passage metadata is retained throughout retrieval and reranking so that each selected passage remains traceable to its original source.


\subsection{Evidence-Sufficiency Gate}
\label{sec:sufficiency-gate}
MedSwin uses an evidence-sufficiency gate to determine whether the retrieved evidence is adequate for answer generation. A set of highly ranked passages may still be insufficient when a required clinical facet is missing, the estimated coverage remains uncertain, or the selected evidence contains unresolved contradictions. The gate evaluates the evidence state before the MedSwin language model is allowed to generate a full response.

\paragraph{Facet-aware evidence-bundle construction}
Let \(\mathcal{C}(q)\) denote the reranked candidate set for clinical query \(q\). MedSwin selects an evidence bundle \(\mathcal{M}\subseteq\mathcal{C}(q)\) subject to both a token budget \(B\) and a maximum number of evidence passages \(K_{\max}\):
\begin{equation}
\label{eq:evidence-bundle-selection}
\small
\begin{aligned}
\mathcal{M}^{\star}
=
\arg\max_{\mathcal{M}\subseteq\mathcal{C}(q)}
&\,U(q,\mathcal{M})\\
\text{s.t.}\quad
&\sum_{d\in\mathcal{M}}\tau(d)\leq B,
\qquad
|\mathcal{M}|\leq K_{\max}.
\end{aligned}
\end{equation}
Here, \(\tau(d)\) is the number of tokens occupied by passage \(d\), measured with the model tokenizer. The token-budget constraint prevents the combined evidence from exceeding the context space allocated to retrieved passages, while \(K_{\max}\) limits the number of chunks included in the bundle.

The utility \(U(q,\mathcal{M})\) favours passages that improve coverage of required clinical facets. It penalises redundant or noisy evidence while preserving safety and contradiction evidence for later adjudication by the sufficiency gate.

MedSwin approximates this constrained objective using greedy selection. For each feasible candidate passage \(d\), it computes the marginal utility
\[
\Delta U(d\mid\mathcal{M})
=
U(q,\mathcal{M}\cup\{d\})-U(q,\mathcal{M}).
\]
The passage with the highest marginal utility per token,
\(\Delta U(d\mid\mathcal{M})/\tau(d)\), is added at each step, provided that its inclusion does not exceed the remaining token budget or the maximum chunk limit. Let \(\pi_j(\mathcal{M})\in[0,1]\) denote the estimated coverage of facet \(f_j\) by bundle \(\mathcal{M}\).

\paragraph{Sufficiency decision}
The gate checks whether every required facet has adequate and sufficiently stable support. A facet is accepted only when its conservative coverage estimate exceeds the configured threshold and its uncertainty remains below the permitted entropy limit. The bundle must also satisfy the contradiction constraints:
\begin{equation}
\label{eq:sufficiency-gate-main}
\small
\begin{aligned}
&\forall f_j\in\mathcal{F}_{\mathrm{critical}}(q):\\
&\quad
\mathrm{LCB}_{1-\delta}
\!\left[\pi_j(\mathcal{M}^{\star})\right]
\geq \theta_j,
\qquad
H_j(\mathcal{M}^{\star})\leq H_{\max},\\
&U_{\mathrm{high}}(\mathcal{M}^{\star})=0,\\
&\left|
\mathcal{C}_{\mathrm{conf}}(\mathcal{M}^{\star})
\right|
\leq c_{\max}.
\end{aligned}
\end{equation}

Here, \(\theta_j\) is the minimum required coverage for facet \(f_j\), and \(H_j\) measures uncertainty in its estimated coverage. \(U_{\mathrm{high}}\) counts unresolved high-severity conflicts, while \(\mathcal{C}_{\mathrm{conf}}\) denotes the set of detected contradictions. The bundle is rejected when a required facet remains insufficient or uncertain, when a high-severity conflict is unresolved, or when the total number of contradictions exceeds the configured limit.

\paragraph{Policy outcomes}
The sufficiency policy produces one of three outcomes:

\begin{enumerate}
    \item \textsc{Answer}: all required facets satisfy the coverage and uncertainty conditions, and the contradiction constraints are met.

    \item \textsc{Retrieve-More}: the gate fails, but another retrieval iteration remains available and additional evidence is expected to improve the unresolved evidence state.

    \item \textsc{Insufficient-Evidence}: the gate fails and further retrieval is unavailable or unlikely to provide sufficient additional utility.
\end{enumerate}

The \textsc{Retrieve-More} decision depends on the remaining retrieval-loop budget and the estimated utility of additional evidence. Missing or uncertain facets determine the target of the next retrieval request. This gate separates MedSwin from a standard top-\(K\) RAG pipeline: retrieval and reranking identify relevant passages, while the sufficiency policy determines whether the resulting evidence state is adequate for generation.

% \subsection{Evidence-Sufficiency Gate - This is needed to be revised}
% \label{subsec:evidence-sufficiency-gate}

% MedSwin uses an \emph{evidence-sufficiency gate} to decide whether the selected evidence bundle is strong enough to support answer generation. In this paper, evidence is sufficient when it covers the clinical facts needed for the question, links guideline or literature evidence to the patient's EMR-derived state, exposes unresolved conflicts, and preserves provenance for review.

% For a clinical query $q$, MedSwin identifies a set of required clinical evidence facets:
% \[
% \mathcal{F}(q)=\{f_1,\ldots,f_m\}.
% \]
% A facet is a decision-relevant information need, such as current medication, recent laboratory status, guideline recommendation, contraindication, dose restriction, patient applicability, or unresolved contradiction.

% An evidence bundle $\mathcal{M}$ is accepted only when it satisfies four conditions:

% \begin{enumerate}
%     \item \textbf{Critical-facet coverage:} all decision-critical facets are supported by accepted evidence.
%     \item \textbf{Patient applicability:} guideline, safety, or literature evidence is connected to the patient's EMR-derived context.
%     \item \textbf{Contradiction handling:} high-severity conflicts are resolved or explicitly preserved as uncertainty.
%     \item \textbf{Traceability:} treatment-relevant claims can be traced to source passages and provenance metadata.
% \end{enumerate}

% Thus, MedSwin does not treat a fixed top-$K$ retrieval set as automatically sufficient. A bundle may contain highly ranked passages but still be insufficient if it misses a contraindication, patient-specific constraint, guideline threshold, or unresolved conflict.

% Operationally, let $\pi_j(\mathcal{M})$ denote the estimated coverage of clinical facet $f_j$ by evidence bundle $\mathcal{M}$. The orchestrator accepts the bundle only if:

% \begin{equation}
% \label{eq:sufficiency-gate-main}
% \small
% \begin{aligned}
% &\min_{j\in\mathcal{F}_{\mathrm{critical}}}
% \mathrm{LCB}_{1-\delta}\!\left[\pi_j(\mathcal{M})\right]
% \ge \theta_j,\\
% &R_{\mathrm{conf}}(\mathcal{M})\le c_{\max},\\
% &T_{\mathrm{trace}}(\mathcal{M})=1.
% \end{aligned}
% \end{equation}

% Here, $\theta_j$ is the required coverage threshold for critical facet $j$, $R_{\mathrm{conf}}$ measures unresolved high-severity contradiction, and $T_{\mathrm{trace}}$ indicates whether the required provenance and citation fields are present.

% The gate produces one of three outcomes:

% \begin{enumerate}
%     \item \textbf{Answer:} the evidence bundle satisfies the gate and the system may generate a grounded answer.
%     \item \textbf{Retrieve more:} one or more critical facets are missing or uncertain, and additional retrieval may improve the evidence state.
%     \item \textbf{Insufficient evidence:} the system cannot satisfy the gate within the available evidence or token budget, so the final response must explicitly state the limitation.
% \end{enumerate}

% This gate is the control point that separates MedSwin from a standard top-$K$ RAG pipeline. Retrieval and reranking construct candidate evidence, but final answer generation is allowed only after the evidence bundle passes the sufficiency gate.

\subsection{Worked Example: Evidence Sufficiency}
\label{subsec:worked-example}

Consider a de-identified medication-safety query:

\begin{quote}
\emph{Can this patient continue metformin after the latest renal-function result?}
\end{quote}

This example illustrates the difference between standard top-$K$ RAG and MedSwin's evidence-sufficiency gate. A standard RAG system may retrieve generally relevant passages about metformin or diabetes management. MedSwin instead checks whether the selected evidence bundle covers the clinical facets required to support the decision: current medication use, latest renal-function status, renal-function trend, relevant guideline threshold, medication-safety warning, and unresolved contradiction.

\begin{table}[h!]
\centering
\caption{Concise evidence-sufficiency check for a medication-safety query.}
\label{tab:worked-example-sufficiency}
\scriptsize
\begin{tabularx}{\linewidth}{X X l X}
\toprule
\textbf{Facet} & \textbf{Accepted evidence} & \textbf{Status} & \textbf{System action} \\
\midrule
Current medication & EMR medication list & Covered & Use in answer \\
Latest renal function & EMR lab result & Covered & Use in answer \\
Renal trend / acute change & Historical labs absent & Missing & Retrieve more or limit answer \\
Guideline threshold & CPG renal-use clause & Covered & Cite guideline \\
Safety warning & Drug-safety passage & Covered & Include caution \\
Contradiction state & No conflict found & Covered & Proceed if missing facet resolved \\
\bottomrule
\end{tabularx}
\end{table}

Because the renal trend is missing, the bundle is not fully evidence-sufficient. MedSwin therefore should not generate a confident recommendation. It must either retrieve additional laboratory history or return a bounded response stating that the available evidence is insufficient to determine whether the latest renal result reflects stable chronic impairment or acute deterioration. This example shows that MedSwin does not accept evidence merely because passages are highly ranked; it accepts evidence only when the required clinical facets are covered and the remaining uncertainty is explicit. A fuller trace of this example is provided in Appendix~\ref{app:worked-example-trace}.

% -
% Figure
% -
\begin{figure*}[t!]
    \centering
    \includegraphics[width=\textwidth]{2sr2.png}
    \caption{Two-stage IR Architecture.}
    \label{fig:ir-pipeline}
\end{figure*}


\subsection{Two-Stage Retrieval-to-Reranker Architecture}
% • Current medical QA models
% • Note: rerankers used in general domain, but NOT finetuned for medical domain
% • Highlight missing work → leads to Contribution (1)
\label{subsec:retrieval-reranking-architecture}

\textbf{Objective.}
The retrieval module has one practical goal: select the smallest set of evidence passages that is strong enough for the clinical question. This evidence may come from the patient's EMR, a CPG document, a drug-safety source, or biomedical literature. A passage is useful only if it helps answer a clinically necessary part of the question.

For example, a medication-safety query may require evidence about: (i) the patient's current medication, (ii) the relevant laboratory result, (iii) the guideline recommendation, (iv) contraindications or dose restrictions, and (v) any unresolved conflict between sources. MedSwin calls these required parts \emph{clinical evidence facets}. The system should not accept an evidence bundle simply because it contains several high-scoring passages. It should accept the bundle only when the critical facets are covered with sufficient confidence.

Formally, for a query $q$, we define a set of clinical facets:
\[
\mathcal{F}(q)=\{f_1,\ldots,f_m\}.
\]
The system then selects an evidence bundle $\mathcal{M}$ under a token budget $B$:
\[
\mathcal{M}^{\star}
=
\arg\max_{\mathcal{M}\subseteq\mathcal{D}}
U(q,\mathcal{M})
\quad
\text{s.t.}
\quad
\sum_{d\in\mathcal{M}}\tau(d)\le B.
\]
Here, $U(q,\mathcal{M})$ rewards evidence that is relevant, reliable, non-redundant, patient-specific, and able to cover the critical facets. It penalises weak provenance, duplicated passages, unsupported claims, and unresolved safety conflicts. This design replaces a fixed top-$K$ context window with a clinically guided selection process.

\subsubsection*{Stage 1: Dense retrieval with vector embedding}~\\
\textbf{Embeddings and similarity.} 
Let $\phi_q,\phi_d:\cdot\rightarrow\mathbb{R}^d$ be the query/document encoders. With L2 normalisation,
\begin{equation}
    \mathbf{q}=\frac{\phi_q(q)}{\|\phi_q(q)\|_2},\quad
    \mathbf{d}=\frac{\phi_d(d)}{\|\phi_d(d)\|_2},\quad
    s_{\mathrm{emb}}(q,d)=\mathbf{q}^{\top}\mathbf{d}.
\end{equation}

\textbf{ANN and lexical union.}
We index $\{\mathbf{d}\}$ with a hybrid ANN (HNSW + IVF-PQ) and retrieve top-$K$ dense candidates $\mathcal{C}_{\mathrm{dense}}(q)$. To handle rare terms/abbreviations, we union with BM25 results $\mathcal{C}_{\mathrm{lex}}(q)$:
\begin{equation}
\label{eq:cand-union}
    \mathcal{C}(q)=\mathrm{TopK}_{K'}\!\big(\mathcal{C}_{\mathrm{dense}}(q)\cup\mathcal{C}_{\mathrm{lex}}(q)\big),\qquad K'\ge K,
\end{equation}
carrying both $s_{\mathrm{emb}}(q,d)$ and $s_{\mathrm{lex}}(q,d)$ for fusion. Time and section filters may be applied using metadata $(t_d,\mathrm{section}_d,\mathrm{source}_d)$.


\subsubsection*{Stage 2: Pointwise LLM reranking}
For each $(q,d)\in\mathcal{C}(q)$, the biomedical reranker from Section~\ref{subsec:biomedical-reranker} produces $\ell_\theta(q,d)$ and the calibrated probability $\widehat{p}_\theta(q,d)$ via \eqref{eq:calib}.


\subsubsection*{Fusion and policy-aware selection}
\textbf{Signal calibration and clinical features.}
For each candidate $(q,d)$, raw dense and lexical similarities are transformed into bounded, query-local calibration features
\[
g_{\mathrm{emb}}(q,d),g_{\mathrm{lex}}(q,d)\in[0,1],
\]
for example using percentile ranks, isotonic calibration, or validation-fitted monotone mappings. This avoids treating embedding cosine scores, BM25 scores, and calibrated reranker probabilities as commensurate quantities.

We augment these retrieval signals with bounded clinical priors:
\begin{equation}
\small
\begin{aligned}
    g_{\mathrm{rec}}(d) &= \exp(-\lambda_r\,\Delta t(d)), \\
    g_{\mathrm{sec}}(d) &\in [0,1] \quad \text{(e.g., CPG recommendation $>$ background)},\\
    g_{\mathrm{src}}(d) &\in [0,1] \quad \text{(source reliability and provenance)}.
\end{aligned}
\end{equation}

\textbf{Fusion score.}
Because dense similarity, lexical similarity, calibrated reranker probability, and clinical priors do not share a common measurement scale, we avoid interpreting their raw weighted sum as a calibrated probability. Instead, each retrieval signal is first mapped to a bounded or calibrated feature. Let
\[
g_{\mathrm{emb}},g_{\mathrm{lex}},g_{\mathrm{rec}},g_{\mathrm{sec}},g_{\mathrm{src}}\in[0,1]
\]
denote percentile- or calibration-transformed versions of the embedding, lexical, recency, section, and source-quality signals. We further define an evidence-quality score
\[
g_{\mathrm{EBM}}(d)\in[0,1],
\]
which encodes the evidence hierarchy of the passage, giving higher weight to current guideline recommendations, systematic reviews, and high-quality randomised trials for generalisable treatment efficacy, while using EMR-derived evidence primarily to assess patient-specific applicability, contraindications, and comorbidity constraints.

The fused score is then computed on the log-odds scale:
\begin{equation}
\label{eq:fusion}
\small
\begin{aligned}
S(q,d)
=
\sigma\big(&
w_0
+\alpha\,\mathrm{logit}(\widehat{p}_\theta(q,d))
+\beta\,g_{\mathrm{emb}}(q,d)
+\gamma\,g_{\mathrm{lex}}(q,d)
\\
&+\rho\,g_{\mathrm{rec}}(d)
+\eta\,g_{\mathrm{sec}}(d)
+\zeta\,g_{\mathrm{src}}(d)
+\xi\,g_{\mathrm{EBM}}(d) 
\\
&-\chi\,g_{\mathrm{noise}}(q,d)
\big),
\end{aligned}
\end{equation}
where $\sigma(\cdot)$ is the logistic function and $g_{\mathrm{noise}}(q,d)$ penalises population mismatch, obsolete recommendations, weak provenance, and semantic drift. This formulation preserves interpretability while preventing unbounded z-scores from being misrepresented as calibrated clinical confidence.

\textbf{Evidence hierarchy and source validity.}
For treatment-oriented decision support, MedSwin assigns each passage an evidence-quality feature $g_{\mathrm{EBM}}(d)$ that favours current guideline recommendations, systematic reviews, and high-quality randomised trials for generalisable efficacy, while using EMR evidence primarily for patient applicability, contraindications, medication constraints, and comorbidity alignment. The full evidence-grade vector and weighting scheme are provided in Appendix~\ref{app:evidence-hierarchy}.


\textbf{Budgeted, diverse selection.}
Let $\tau(d)$ denote the token cost of passage $d$. Rather than enforcing raw source-count constraints, we select an evidence bundle by maximising clinical utility under a hard token budget:
\begin{equation}
\label{eq:budgeted-utility-selection}
\small
\mathcal{M}^{\star}
=
\arg\max_{\mathcal{M}\subseteq\mathcal{C}(q)}
U(q,\mathcal{M})
\quad
\text{s.t.}
\quad
\sum_{d\in\mathcal{M}}\tau(d)\le B.
\end{equation}
Diversity is incorporated through the redundancy penalty $R_{\mathrm{red}}(\mathcal{M})$ inside $U(q,\mathcal{M})$, rather than through a separate MMR objective. This avoids maintaining two competing optimisation criteria and allows redundancy, contradiction, evidence quality, and safety coverage to be adjudicated within a single clinically interpretable utility function.

\textbf{Greedy solver, facet sufficiency, and stopping.}
We approximate the constrained optimisation with a greedy marginal-utility scheme. For each candidate passage $d$ and clinical facet $f_j\in\mathcal{F}(q)$, let
\[
a_{dj}=P(f_j \text{ is adequately addressed by } d\mid q)
\]
be a calibrated facet-alignment score produced by the reranker or a lightweight verifier. The probability that the current bundle $\mathcal{M}$ covers facet $f_j$ is modelled as
\begin{equation}
\label{eq:facet-coverage}
\small
\pi_j(\mathcal{M})
=
1-\prod_{d\in\mathcal{M}}
\left(
1-a_{dj}\widehat{p}_\theta(q,d)g_{\mathrm{EBM}}(d)
\right).
\end{equation}
Because retrieved passages are not necessarily independent, near-duplicate passages or passages derived from the same parent source are down-weighted through $R_{\mathrm{red}}(\mathcal{M})$ and grouped by provenance before computing effective coverage. Thus, corroboration is credited primarily when it arises from independent sources or distinct evidence strata.
This noisy-OR form rewards independent corroboration while saturating once additional passages become redundant.

The clinical utility of a bundle is
\begin{equation}
\label{eq:clinical-utility}
\tiny
U(q,\mathcal{M}) = \sum_{j=1}^{m} v_j \pi_j(\mathcal{M}) + \lambda_{\mathrm{safety}} R_{\mathrm{safety}}(\mathcal{M}) - \lambda_{\mathrm{red}} R_{\mathrm{red}}(\mathcal{M}) - \lambda_{\mathrm{conf}} R_{\mathrm{conf}}(\mathcal{M}),
\end{equation}
where $v_j$ is the clinical importance of facet $f_j$, $R_{\mathrm{safety}}$ rewards coverage of contraindications and high-severity risks, $R_{\mathrm{red}}$ penalises near-duplicate evidence, and $R_{\mathrm{conf}}$ penalises unresolved contradictions among high-quality sources.

At each step, the system selects the feasible passage with the highest positive marginal gain 
\begin{equation}
\label{eq:marginal-utility}
\Delta U(d\mid\mathcal{M})
=
U(q,\mathcal{M}\cup\{d\})-U(q,\mathcal{M}),
\end{equation}
subject to the remaining token budget. The system does not stop merely because top-level CPG and EMR counts are satisfied. Instead, the orchestrator accepts the evidence bundle only when all of the following conditions hold:
\begin{equation}
\label{eq:stopping}
\begin{aligned}
&\min_{j\in\mathcal{F}_{\mathrm{critical}}}
\mathrm{LCB}_{1-\delta}\!\left[\pi_j(\mathcal{M})\right]
\ge \theta_j,
\\
&\max_{j\in\mathcal{F}_{\mathrm{critical}}}
H\!\left(\pi_j(\mathcal{M})\right)
\le h_{\max},
\\
&R_{\mathrm{conf}}(\mathcal{M})\le c_{\max},
\\
&\max_{d\in\mathcal{C}(q)\setminus\mathcal{M}}
\frac{\Delta U(d\mid\mathcal{M})}{\tau(d)}
\le \epsilon_U.
\end{aligned}
\end{equation}
Here $\mathrm{LCB}_{1-\delta}$ is a lower confidence bound on facet coverage, $H(\cdot)$ is binary entropy, $\theta_j$ is the required clinical sufficiency threshold for facet $j$, and $\epsilon_U$ is the minimum marginal utility per token required to justify further retrieval. If any critical facet remains uncertain, if high-quality sources conflict, or if the expected marginal utility of additional retrieval exceeds $\epsilon_U$, the orchestrator issues a targeted \textit{retrieve-more} action. This action may increase $K'$, expand synonyms, relax metadata filters, or dispatch a specialist agent to search for missing facets such as contraindications, population exclusions, or rare adverse events.

\textbf{Implementation details.}
After evidence selection, MedSwin applies facet-aware context packing to remove redundant or low-utility passages while preserving contraindications, subgroup exclusions, and high-grade contradictory evidence. LCB estimation, protected-pruning rules, chunk metadata, ANN settings, and runtime complexity are provided in Appendix~\ref{app:retrieval-implementation}.

\subsection{Multi-Agent Coordination}
\label{subsec:multi-agent-coordination}

The Multi-Agent Coordination (MAC) layer sits above retrieval and reranking. Its role is not to make the system more complex or to let agents freely deliberate. Instead, MAC performs targeted evidence exploration for missing clinical facets before final synthesis. Each agent searches for a different type of evidence, such as EMR patient context, guideline recommendations, contraindications, safety warnings, evidence quality, or contradictions. The orchestrator then applies the same evidence-sufficiency gate described in Section~\ref{sec:evidence-sufficiency}.

\begin{figure*}[t]
    \centering
    \includegraphics[width=\textwidth]{Multi-Agent.png}
    \caption{Multi-agent coordination for evidence-gated clinical RAG. Agents retrieve and verify facet-specific evidence; the orchestrator accepts, expands, or limits the evidence bundle before final answer generation.}
    \label{fig:multi-agent-coordination}
\end{figure*}

MAC uses six bounded agents. The \textbf{EMR agent} retrieves patient-specific facts such as medications, laboratory values, comorbidities, allergies, and recent events. The \textbf{guideline agent} retrieves CPG recommendations, eligible populations, contraindicated populations, recommendation strength, and guideline version. The \textbf{safety agent} searches for contraindications, adverse events, dose restrictions, drug-drug interactions, and high-severity warnings. The \textbf{evidence-quality agent} assigns source type, evidence grade, recency, population fit, and provenance quality. The \textbf{contradiction critic} identifies conflicts between high-quality sources and distinguishes genuine disagreement from outdated evidence or population mismatch. The \textbf{synthesis agent} generates the final response only from the accepted evidence bundle and sufficiency state.

Each agent emits structured claims rather than free-text summaries:
\[
c =
(f_j,\mathrm{claim},\mathrm{polarity},d,g_{\mathrm{EBM}},
\widehat{p}_{\theta},\mathrm{provenance}),
\]
where $f_j$ is the clinical facet, $d$ is the supporting passage, $\mathrm{polarity}$ indicates whether the passage supports, contradicts, or qualifies the claim, $g_{\mathrm{EBM}}$ is the evidence-quality score, and $\widehat{p}_{\theta}$ is the calibrated reranker relevance score. This structure prevents the final synthesiser from treating retrieved context as an undifferentiated text block.

If the sufficiency gate is not satisfied, MAC issues a targeted \textit{retrieve-more} action rather than simply increasing top-$K$. Missing contraindication evidence triggers the safety agent; weak guideline coverage triggers the guideline agent; missing patient applicability triggers the EMR agent; and unresolved conflicts trigger the contradiction critic. If the missing facet cannot be resolved, the final answer must explicitly state that the evidence is insufficient.

For each query, MAC returns four audit artefacts: a retrieval trace, a facet-coverage matrix, a contradiction ledger, and a final sufficiency decision. These artefacts show whether MedSwin answered because the evidence was sufficient, retrieved more because a facet was missing, or produced a limited response because the available evidence did not support a full recommendation. Full MAC aggregation, reliability weighting, contradiction adjudication, and audit-record details are provided in Appendix~\ref{app:agent-artifacts}.


\section{Experiments}
\label{sec:eval_medquad}

\subsection{Evaluation Design}
\subsubsection{Standalone Medical LLM Evaluation}
We first evaluate the answer-generation capability of MedSwin independently from the retrieval and multi-agent orchestration components. This evaluation isolates whether the proposed student models, including supervised fine-tuning, knowledge distillation, sequential fine-tuning, and training-free merged variants, improve biomedical answer generation relative to the base MedAlpaca-7B model, other open-source biomedical baselines, and frontier general-purpose models. We evaluate on MedQuAD \cite{ben2019question} and HealthBench \cite{arora2025healthbench}, which provide complementary settings for medical question answering and health-related response generation. The standalone evaluation is not intended to establish clinical safety. Instead, it measures whether the MedSwin language model variants improve semantic alignment, lexical faithfulness, and response precision under controlled benchmark conditions. This separation is important because a model may produce semantically plausible answers while still failing to retrieve patient-specific evidence or expose sufficient clinical provenance.

\subsubsection{Biomedical Reranker Evaluation}
Second, we evaluate MedSwin-Reranker as an independent evidence-ranking component. The goal of this experiment is to determine whether biomedical adaptation improves the ordering of candidate evidence passages before they are passed into the system-level evidence sufficiency module. The reranker is evaluated under a fixed candidate-generation setting, where each query is paired with a pool of candidate passages and the reranker must assign higher scores to clinically or scientifically relevant evidence. 

\textbf{still on progress for this part}

\subsubsection{End-to-End System Evaluation}
Third, we evaluate the complete MedSwin runtime as an evidence-grounded clinical decision-support system. This experiment tests the full pipeline, including retrieval, reranking, evidence sufficiency gating, multi-agent verification, citation production, and final answer synthesis. We use TREC Clinical Decision Support 2016 as a system-level audit benchmark because each case requires patient-context reasoning and evidence retrieval from a biomedical corpus. Unlike standalone medical QA evaluation, the system-level evaluation focuses on whether MedSwin can retrieve sufficient evidence, identify critical clinical facets, avoid unsupported recommendations, and expose an auditable trace. This design reflects the central claim of MedSwin that clinical decision-support reliability depends on both fluent answer generation and on evidence sufficiency, provenance, and controlled abstention when the available evidence is incomplete.

\subsection{Metrics}
\subsubsection{Generation Metrics}
For standalone medical LLM evaluation, we report ROUGE-L F1, BERTScore-F1, token-level F1, and unigram/bigram precision. ROUGE-L and token-level F1 measure lexical overlap, while BERTScore-F1 captures contextual semantic similarity \cite{lin2004rouge,bertscore}. Unigram and bigram precision are used as reference-precision proxies, indicating whether responses remain concise and close to the reference. For visualization only, we also group these metrics into an Answer Quality Index and a Reference Precision Index. The raw metrics remain the primary results; the indices are used only to make multi-metric trade-offs, following standard multi-metric benchmark practice \cite{wang2018glue,wang2019superglue,liang2023holistic,muennighoff2023mteb}.


\begin{table*}[h!]
    \centering
    % LEFT TABLE (MedQuAD)
    \begin{minipage}[h]{0.48\textwidth}
        \centering
        \caption{MedQuAD benchmarking results.}
        \label{tab:medquad_bench}
        \scriptsize
        \begin{tabular}{lccccc}
        \toprule
        \textbf{Model} & \textbf{rougeL\_f} & \textbf{bert\_f} & \textbf{tok\_f1} & \textbf{uni\_prec} & \textbf{bi\_prec} \\
        \midrule
        DaRE-TIES-KD-0.7       & 0.17 & 0.8501 & 0.22 & 0.47 & 0.14 \\
        NuSL-KD-0.7            & 0.17 & 0.8500 & 0.21 & 0.47 & 0.11 \\
        DaRE-Linear-KD-0.7     & 0.17 & 0.8494 & 0.22 & 0.45 & 0.14 \\
        TA-SFT-0.7             & 0.16 & 0.8491 & 0.18 & 0.51 & 0.20 \\
        7B-KD                  & 0.18 & 0.8441 & 0.26 & 0.46 & 0.12 \\
        7B-SFT                 & 0.15 & 0.8396 & 0.18 & 0.52 & 0.21 \\
        KD-SFT-PubMed-l        & 0.17 & 0.8341 & 0.24 & 0.44 & 0.11 \\
        KD-SFT-PubMed-map      & 0.17 & 0.8297 & 0.22 & 0.47 & 0.12 \\
        MedGemma-27b-it        & 0.19 & 0.8465 & 0.27 & 0.40 & 0.09 \\
        MedGemma-1.5-4b        & 0.18 & 0.8449 & 0.26 & 0.40 & 0.08 \\
        Gemini 3.1             & 0.15 & 0.8490 & 0.26 & 0.36 & 0.08 \\
        Gemini 2.5 Flash       & 0.09 & 0.8390 & 0.11 & 0.55 & 0.28 \\
        GPT 5.4                & 0.13 & 0.8450 & 0.21 & 0.31 & 0.05 \\
        Grok 4.2 (NR)          & 0.13 & 0.8333 & 0.22 & 0.29 & 0.05 \\
        Grok 4.1 (NR)          & 0.12 & 0.8340 & 0.18 & 0.28 & 0.05 \\
        Claude Sonnet 4.6      & 0.12 & 0.8470 & 0.19 & 0.28 & 0.05 \\
        Mistral Small          & 0.14 & 0.8450 & 0.23 & 0.34 & 0.07 \\
        Mistral Large          & 0.13 & 0.8430 & 0.21 & 0.29 & 0.06 \\
        BioMistral-7B          & 0.18 & 0.8460 & 0.25 & 0.45 & 0.11 \\
        Meditron-7B            & 0.16 & 0.8240 & 0.22 & 0.47 & 0.13 \\
        MedAlpaca-7B           & 0.16 & 0.8196 & 0.18 & 0.50 & 0.20 \\
        \bottomrule
        \end{tabular}
    \end{minipage}\hfill
    % RIGHT TABLE (HealthBench)
    \begin{minipage}[h]{0.48\textwidth}
        \centering
        \caption{OpenAI HealthBench benchmarking results.}
        \label{tab:healthbench_bench}
        \scriptsize 
        \begin{tabular}{lccccc}
        \toprule
        \textbf{Model} & \textbf{rougeL\_f} & \textbf{bert\_f} & \textbf{tok\_f1} & \textbf{uni\_prec} & \textbf{bi\_prec} \\
        \midrule
        % TIES-KD-0.7-0.6            & 0.16 & 0.8411 & 0.22 & 0.35 & 0.08 \\
        TA-SFT-0.7                 & 0.16 & 0.8410 & 0.23 & 0.35 & 0.08 \\
        NuSL-KD-0.7                & 0.16 & 0.8403 & 0.24 & 0.34 & 0.07 \\
        DaRE-TIES-KD-0.7           & 0.16 & 0.8392 & 0.24 & 0.34 & 0.07 \\
        DaRE-Linear-KD-0.7         & 0.16 & 0.8388 & 0.24 & 0.34 & 0.07 \\
        7B-KD                      & 0.15 & 0.8314 & 0.23 & 0.32 & 0.06 \\
        7B-SFT                     & 0.13 & 0.8299 & 0.17 & 0.31 & 0.07 \\
        KD-SFT-PubMed-l            & 0.13 & 0.8186 & 0.20 & 0.32 & 0.06 \\
        KD-SFT-PubMed-map          & 0.12 & 0.8088 & 0.16 & 0.31 & 0.06 \\
        MedGemma-27b-it            & 0.17 & 0.8342 & 0.26 & 0.32 & 0.07 \\
        MedGemma-1.5-4b            & 0.16 & 0.8326 & 0.25 & 0.32 & 0.06 \\
        Gemini 3.1                 & 0.14 & 0.8480 & 0.26 & 0.33 & 0.06 \\
        Gemini 2.5 Flash           & 0.08 & 0.8230 & 0.13 & 0.47 & 0.18 \\
        GPT 5.4                    & 0.12 & 0.8460 & 0.22 & 0.27 & 0.04 \\
        Grok 4.2 (NR)              & 0.12 & 0.8390 & 0.23 & 0.27 & 0.04 \\
        Grok 4.1 (NR)              & 0.11 & 0.8310 & 0.18 & 0.26 & 0.04 \\
        Claude Sonnet 4.6          & 0.11 & 0.8440 & 0.20 & 0.24 & 0.04 \\
        Mistral Small              & 0.13 & 0.8470 & 0.23 & 0.30 & 0.06 \\
        Mistral Large              & 0.12 & 0.8400 & 0.21 & 0.26 & 0.04 \\
        BioMistral-7B              & 0.16 & 0.8404 & 0.23 & 0.35 & 0.08 \\
        Meditron-7B                & 0.12 & 0.8068 & 0.15 & 0.31 & 0.06 \\
        MedAlpaca-7B               & 0.11 & 0.7946 & 0.21 & 0.34 & 0.07 \\
        \bottomrule
        \end{tabular}
    \end{minipage}
\end{table*}

\begin{figure*}[h!]
    \centering
    % LEFT FIGURE (MedQuAD)
    \begin{minipage}[h]{0.48\textwidth}
        \centering
        \includegraphics[width=\linewidth]{medquad_dashboard.png}
        \caption{MedQuAD evaluation over Answer Quality Index and Grounding Precision Index.}
        \label{fig:quad_placeholder}
    \end{minipage}\hfill
    % RIGHT FIGURE (HealthBench)
    \begin{minipage}[h]{0.48\textwidth}
        \centering
        \includegraphics[width=\linewidth]{healthbench_dashboard.png}
        \caption{OpenAI HealthBench evaluation over Answer Quality Index and Grounding Precision Index.}
        \label{fig:health_placeholder}
    \end{minipage}
\end{figure*}


\subsubsection{Reranker Metrics}
For biomedical reranker evaluation, we will report MRR@10, nDCG@10, Recall@10, MAP@10, and p95 latency. MRR@10 measures whether the first relevant evidence passage is ranked near the top, while nDCG@10 captures graded ranking quality across the top results. Recall@10 measures whether relevant passages are retained within the context budget, and MAP@10 summarizes precision across ranked positions. We additionally report p95 latency to assess whether the reranker is practical for deployment in an interactive clinical decision-support workflow.

\textbf{still on going}

\subsubsection{System-Level Clinical Audit Metrics}
For the end-to-end system audit, we report facet recall (FR), critical facet recall (CFR), evidence-document recall (EDR), citation precision, groundedness proxy (GP), trace completeness (TC), sufficiency decision score (SDS), unsupported penalty (USP), unsafe omission penalty (UOP), policy pass rate, degraded rate, and error rate. These metrics are designed to evaluate whether the system retrieves clinically necessary evidence, grounds its answer in the retrieved record, and correctly identifies cases where evidence is insufficient for a full recommendation.

We also define the composite MedSwin System Audit Score (MSAS), which rewards evidence sufficiency and traceable grounding while penalising unsupported or safety-critical omissions:
\begin{equation}
\label{eq:msas}
\scriptsize
\begin{aligned}
\text{MSAS} = \mathrm{clip}_{[0,1]}\Big(&
0.25 \cdot \text{CFR}
+0.15 \cdot \text{FR}
+0.15 \cdot \text{GP}
+0.15 \cdot \text{SDS}
+0.10 \cdot \text{TC}
\\
&+0.10 \cdot \text{EDR}
+0.10 \cdot \text{CQP}
-0.20 \cdot \text{UOP}
-0.10 \cdot \text{USP}
\Big),
\end{aligned}
\end{equation}
where the clipping operation constrains the final score to $[0,1]$. The weighting deliberately gives highest positive weight to critical-facet recall because, in clinical decision support, omission of a contraindication, patient-applicability constraint, or treatment-relevant exclusion can be more consequential than minor lexical mismatch with a reference answer. MSAS is therefore not intended to replace classical retrieval or generation metrics. Instead, it operationalises the specific safety contract claimed by MedSwin that retrieve the right evidence, know when evidence is insufficient, ground the answer, and expose the trace.


\paragraph{Linking the worked example to evaluation}
The worked example in Section~\ref{subsec:worked-example} illustrates the behaviour that the system-level audit is designed to measure. A successful system should not merely generate a fluent answer. It should retrieve the necessary EMR and CPG evidence, rank clinically relevant passages above background information, identify whether critical facets are covered, and expose the provenance of the final answer. The system-level metrics therefore evaluate the same behaviours shown in the example: critical-facet recall, evidence-document recall, citation precision, groundedness, trace completeness, and sufficiency-decision quality.


\section{Results}

\subsection{System evaluation results}
Table~\ref{tab:trec_cds_system_audit} shows that MedSwin's reliability emerges from the interaction between retrieval, reranking, facet sufficiency, and evidence-gated generation. The MedSwin backend obtains the highest policy pass rate, critical-facet recall, evidence-document recall, trace completeness, and tied-best MSAS, with zero runtime audit errors. This indicates that the MedSwin-specific model stack better preserves clinically pivotal evidence under the same system architecture. Detailed metric-by-metric interpretation is provided in Appendix~\ref{app:detailed-system-results}.

\paragraph{Clinical decision-support reliability}
The whole-system audit supports the central claim that MedSwin improves clinical decision-support reliability by making evidence sufficiency explicit. In treatment and diagnostic settings, the risk is often not that a model fails to produce a fluent answer, but that it produces a plausible answer while omitting a critical contraindication, exclusion criterion, patient-specific constraint, or unresolved evidence conflict. MedSwin addresses this failure mode by requiring critical facets to be covered before answer generation is accepted. The observed critical-facet recall of 0.82 under the MedSwin backend demonstrates that the architecture is effective at surfacing pivotal information, while the policy pass rate of 0.85 shows that the system can translate evidence adequacy into generation decisions rather than answering indiscriminately.

This distinction is essential for clinical AI evaluation. A generic medical LLM can score well on semantic similarity while still being unsafe if its answer is unsupported by the retrieved record. Conversely, a cautious system may produce fewer full recommendations but provide safer bounded responses when evidence is incomplete. MedSwin is evaluated under the latter standard: correctness is not defined only by textual similarity to a reference answer, but by whether the system retrieves sufficient evidence, identifies missing facets, exposes provenance, and prevents unsupported recommendations. This makes the benchmark more demanding than answer-only QA and better aligned with clinical decision-support deployment.


\subsection{Results and interpretation}
The results show a consistent separation between lexical-match strength and semantic alignment, and make the value of merging explicit across two benchmarks. 

Figures~\ref{fig:quad_placeholder} and  ~\ref{fig:health_placeholder} summarize the standalone benchmark results using two visualization-only indices. The scatter plots separate semantic/lexical answer quality from reference precision, while the ranked panels provide compact model ordering. The following interpretation relies on the raw values in Tables~\ref{tab:medquad_bench} and~\ref{tab:healthbench_bench}.

On MedQuAD, KD alone remains the strongest \emph{student training} baseline on strict lexical proxies when achieving $0.26$ at token F1 and $0.18$ ROUGE-L (Table \ref{tab:medquad_bench}), consistent with teacher imitation improving fine-grained token selection and reference coverage. SFT alone yields the highest unigram/bigram precision with uni\_prec=$0.52$; bi\_prec=$0.21$ (Table \ref{tab:medquad_bench}), consistent with instruction tuning bias toward more conservative, templated outputs that reduce over-generation.

Crucially, \textbf{training-free merging yields the best semantic alignment at fixed model size} on MedQuAD: DARE-TIES attains the highest BERTScore-F1 ($0.8501$), narrowly ahead of NuSL ($0.8500$), DARE-Linear ($0.8494$), and TA ($0.8491$). This improves over MedAlpaca-7B ($0.8196$, +$0.0305$) and also exceeds KD-only ($0.8441$) and SFT-only ($0.8396$). The frontier baselines are competitive on semantics, with Gemini 3.1 ($0.8490$), Claude Sonnet 4.6 ($0.8470$), GPT-5.4 ($0.8450$), and Mistral Large ($0.8430$) clustering near the top, but MedSwin still holds the best semantic score while keeping stronger precision. This indicates that merging primarily closes the semantic-alignment gap without necessarily optimizing lexical exact-match proxies.

On HealthBench, the same pattern largely persists under a broader comparison set. The best semantic result is obtained by TIES-KD-0.7-0.6 (BERTScore-F1 $0.8411$), followed extremely closely by TA-SFT-0.7 ($0.8410$), BioMistral-7B ($0.8404$), NuSL-KD-0.7 ($0.8403$), and DARE-TIES-KD-0.7 ($0.8392$). The frontier baselines again sit in the same semantic band, with Gemini 3.1 ($0.8480$), GPT-5.4 ($0.8460$), Claude Sonnet 4.6 ($0.8440$), Mistral Small ($0.8470$), and Mistral Large ($0.8400$) showing strong semantic alignment but weaker precision than the best MedSwin variants. In contrast, the MedGemma baselines remain strongest on lexical overlap proxies, with MedGemma-27B achieving the top ROUGE-L/token-F1 scores on HealthBench ($0.1720/0.2574$), followed by MedGemma-1.5 ($0.1629/0.2523$). Relative to the base MedAlpaca-7B, however, every major MedSwin variant improves substantially on HealthBench semantics, with the best merge increasing BERTScore from $0.7946$ to $0.8411$ (+$0.0465$). The cross-benchmark consistency is important: although the winning operator changes slightly by dataset and metric, training-free composition repeatedly delivers the best or near-best semantic alignment among 7B-scale students, while KD preserves lexical faithfulness and SFT improves precision-oriented response restraint.

\begin{table}[h!]
\centering
\caption{MedSwin System evaluation on TREC Clinical Decision Support 2016.}
\label{tab:trec_cds_system_audit}
\scriptsize
\begin{tabular}{lcccccccc}
\toprule
\textbf{Ecosystem} 
& \textbf{PPR} 
& \textbf{FR} 
& \textbf{CFR} 
& \textbf{EDR} 
& \textbf{CitPrec} 
& \textbf{TC} 
& \textbf{MSAS} 
& \textbf{Err.} \\
\midrule
OpenAI\footnotemark
& 0.78 & \textbf{0.92} & 0.78 & 0.45 & \textbf{0.88} & 0.78 & 0.84 & 0.00 \\
GoogleAI\footnotemark
& 0.76 & 0.88 & 0.72 & 0.50 & 0.81 & 0.76 & 0.75 & 0.00 \\
MedSwin\footnotemark
& \textbf{0.85} & 0.90 & \textbf{0.82} & \textbf{0.53} & 0.77 & \textbf{0.81} & \textbf{0.84} & 0.00 \\
\bottomrule
\end{tabular}
\vspace{-0.5em}
\end{table}

\footnotetext{OpenAI Backend: gemini-embedding-2 + gemini-3.0}
\stepcounter{footnote}
\footnotetext{Google Backend: text-embedding-3-large + gpt-5.4}
\stepcounter{footnote}
\footnotetext{MedSwin: DaRE-TIES-KD-0.7 + medswin-reranker}

The system audit further clarifies the role of the MedSwin models under architectural constraints. Standalone model benchmarks show that training-free merges improve semantic alignment at 7B scale, while reranker benchmarks show that MedSwin-Rerank provides the strongest ranking quality among the compared rerankers. The TREC-CDS audit tests whether these improvements remain useful when the models are embedded inside a full clinical RAG pipeline. The MedSwin backend's leading critical-facet recall and evidence-document recall indicate that the specialised student and reranker are not merely optimising benchmark overlap; they improve the retrieval and preservation of clinically relevant evidence under token and the evidence-sufficiency gate. This is a stronger reliability claim than standalone answer quality: it shows that the MedSwin model stack supports the architecture's evidence-selection objective rather than only producing more semantically similar text.

This pattern is consistent with interference-aware composition: DARE reduces dense overlap among potentially conflicting coordinates, while TIES-style sign resolution prevents remaining overlap from canceling on high-signal parameters \cite{yu2023supermario,yadav2023ties}. The small but consistent advantage of DARE-TIES over DARE-Linear on MedQuAD, together with the strong HealthBench result of the TIES-based merge, suggests that explicit conflict handling contributes materially when composing SFT and KD updates.


\paragraph{Why sequential SFT-on-KD degrades semantic alignment}
Sequential SFT-on-KD performs worse than training-free merging on semantic-alignment metrics, suggesting that continued fine-tuning can partially overwrite teacher-like behaviour learned during distillation. A detailed analysis is provided in Appendix~\ref{app:sequential-sft-kd}.

\paragraph{Evaluation caveats and impact}
Overlap and semantic-similarity metrics remain useful for scalable model comparison, but they do not establish clinical safety. A response can obtain a high BERTScore while omitting a contraindication, citing irrelevant evidence, or failing to disclose uncertainty. We therefore treat ROUGE-L, token-F1, BERTScore-F1, and n-gram precision as model-level indicators rather than system-level safety measures. The system audit closes this gap by evaluating evidence coverage, sufficiency decisions, citation precision, and traceability under a patient-context retrieval setting. This distinction is central to MedSwin's contribution: the architecture is not optimised merely to generate medically plausible text, but to decide when a grounded answer is warranted and when the evidence remains insufficient.

The results suggest two deployment-relevant conclusions. First, training-free merging provides a practical path to stronger 7B-scale medical generation without repeated fine-tuning or ensembling, which is important for resource-constrained clinical assistants \cite{wortsman2022soups,goddard2024arcee}. Second, the MedSwin backend remains competitive with frontier-backed versions of the same architecture in the full system audit, tying the best MSAS while achieving the highest critical-facet recall, evidence-document recall, policy pass rate, and trace completeness. This supports the feasibility of a compact, auditable medical AI stack in settings where external frontier APIs may be undesirable due to cost, governance, privacy, or deployment constraints.

\begin{table}[h!]
\centering
\caption{MedQuAD retrieval-style reranker benchmarking results.}
\label{tab:medquad_rerank_bench}
\footnotesize 
\setlength{\tabcolsep}{1.5pt} 
\begin{tabular}{lccccc}
\toprule
Model & MRR@10 & nDCG@10 & Recall@10 & MAP@10 & p95 ms \\
\midrule
MedSwin-Rerank & 0.9771 & 0.9826 & 1.0000 & 0.9771 & 3895.1 \\
BGE-v2-Gemma & 0.9733 & 0.9799 & 0.9990 & 0.9733 & 3670.5 \\
Qwen3-Reranker-8B$^\dagger$ & 0.9739 & 0.9798 & 0.9977 & 0.9739 & \textit{rerun} \\
Jina-v3-Rerank$^\dagger$ & 0.9700 & 0.9768 & 0.9973 & 0.9700 & \textit{rerun} \\
BGE-v2-M3 & 0.9684 & 0.9755 & 0.9972 & 0.9684 & 209.4 \\
\bottomrule
\end{tabular}
\end{table}

\paragraph{Reranker benchmarking on MedQuAD (retrieval-style)}
% TO BE UPDATED ONCE MTEB BENCHMARK IS DONE
In addition to answer-generation scoring, we benchmark \textbf{MedSwin-Rerank} on MedQuAD under a retrieval-style formulation and report MRR@10, nDCG@10, Recall@10, MAP@10, and p95 latency. Candidate construction and scoring are described in Appendix~\ref{sec:appendix}. MedSwin-Rerank achieves the strongest ranking quality among the compared rerankers, with MRR@10 of 0.9771, nDCG@10 of 0.9826, and Recall@10 of 1.0000. BGE-v2-Gemma remains close, indicating that the base reranker is already strong, while the MedSwin adaptation provides a measurable improvement in ranking quality. BGE-v2-M3 is substantially faster but slightly weaker, making it a useful efficiency baseline rather than the strongest evidence-selection model.


\section{Discussion and Limitations}
\label{sec:discussion-limitations}

\subsection{What the results show}
The main result is that MedSwin improves evidence-gated clinical RAG behaviour, not merely answer fluency. The TREC-CDS audit shows that MedSwin retrieves and preserves more critical evidence under the same system contract, with higher critical-facet recall, evidence-document recall, policy pass rate, and trace completeness. This supports the paper's central claim: clinical RAG should be evaluated by whether it retrieves sufficient evidence and exposes a reviewable trace, not only by whether the final answer is semantically similar to a reference.

The reranker and model results explain why the system behaves this way. The biomedical reranker improves the ordering of candidate evidence, and the compact MedSwin LLM provides a deployable generator that remains competitive on standalone medical QA benchmarks. However, neither component alone is the novelty. The novelty is the sufficiency-controlled interaction among retrieval, reranking, agentic verification, and bounded generation.

\subsection{Why this is different from generic RAG}
Generic RAG asks: \emph{What passages should be placed in the prompt?} MedSwin asks: \emph{Is the available evidence enough to support this clinical answer?} This changes both system behaviour and evaluation. A high-ranked passage is not automatically accepted; it must support a required clinical facet. A fluent answer is not automatically valid; it must be grounded in the accepted evidence bundle. A missing contraindication or patient-applicability constraint is treated as a system failure even if the final response sounds medically plausible.

\subsection{Limitations}
MedSwin is not a clinically validated autonomous decision-maker. The TREC-CDS audit evaluates evidence retrieval, sufficiency behaviour, and traceability under a reproducible benchmark, but it does not replace prospective clinician evaluation or local safety validation. TREC relevance judgments are document-level rather than claim-level clinical adjudications, and the seeded facets cannot capture every nuance of bedside decision-making.

The sufficiency thresholds are operational policy parameters. They make the system auditable, but they require local validation before deployment. Different institutions may require different thresholds for different query types, specialties, patient populations, and risk levels. Future work should study clinician-annotated facet labels, prospective safety review, and threshold calibration under real institutional governance.

Finally, MedSwin improves evidence control but does not eliminate medical AI risk. Retrieval may miss relevant evidence, guidelines may be outdated, EMR data may be incomplete, and source conflicts may require human adjudication. For this reason, MedSwin is best understood as clinician-facing decision support with explicit evidence boundaries, not as a replacement for clinical judgment.


\section{Open-Source and Reproducibility}
% • Current lack of open medical multi-agent frameworks
% • Community demand for transparent, reproducible systems
% • Highlight missing work → leads to Contribution (3)
\label{subsec:opensource-reproducibility}

Reproducibility is a prerequisite for biomedical translation: without verifiable pipelines, improvements in medical QA can reflect unreported degrees of freedom (data mixture, augmentation filters, teacher decoding, retrieval policies, or evaluation prompts) rather than genuine methodological progress. This concern is emphasized repeatedly in medical AI and clinical ML guidance, which calls for transparent reporting, traceable artifacts, and reproducible experimentation before deployment claims are credible \cite{vollmer2020mlhealth,beam2018mlhealth}. MedSwin is therefore designed as a \emph{reference implementation} with end-to-end scripting, provenance tracking, and release artifacts aligned with open science principles.

\subsection{Release scope and artifacts}
We will release a complete reproducibility package containing:
\begin{itemize}
    \item \textbf{Data engineering and safety filters:} license-aware dataset fetchers; preprocessing; PHI/identifier scrubbing; MD5-based deduplication; medically-invariant augmentation filters; and dataset manifests capturing transformation provenance per sample.
    \item \textbf{SFT and KD training code:} QLoRA/LoRA configurations, optimizer and schedule settings, gradient checkpointing knobs, and training scripts that reconstruct MedSwin-SFT and MedSwin-KD checkpoints from raw inputs \cite{hu2021lora,dettmers2023qlora}.
    \item \textbf{Teacher-label generation utilities:} fully logged prompts and decoding parameters for hard-label completion, plus serialization utilities for top-$k$ soft-label distillation traces sourced from MedGemma-27B-it \cite{google2025medgemma,hinton2015distilling}.
    \item \textbf{Model merging specifications:} MergeKit YAML configurations and reproducible scripts for TA/SLERP/DARE/TIES compositions, including merge coefficients used for reported results \cite{goddard2024arcee,yadav2023ties,yu2023supermario,shoemake1985quaternion}.
    \item \textbf{Retrieval + reranking pipeline:} scripts to reproduce the two-stage IR pipeline (dense + lexical union, long-context biomedical reranker inference, calibrated scoring hooks), implemented using open reranker tooling and documented index parameters \cite{lewis2020rag,bgev2gemma,flagembedding}.
    \item \textbf{Evaluation and reporting:} deterministic prompting templates, decoding parameters, and metric computation for MedQuAD, HealthBench, and the system audit benchmark; plus optional RAG evaluation hooks to assess retrieval quality and answer groundedness under a standardized interface.
\end{itemize}

\subsection{FAIR-aligned provenance and governance}
We operationalize open-science principles by capturing dataset provenance and experimental metadata in machine-readable form, consistent with the FAIR emphasis on findability, accessibility, interoperability, and reusability \cite{wilkinson2016fair}. Each training/evaluation artifact is accompanied by:
\begin{itemize}
    \item dataset hashes, filtering logs, and augmentation lineage;
    \item checkpoint lineage (base model, delta source, merge operator, coefficients);
    \item evaluation manifests (prompt template version, decoding params, metrics, and run identifiers).
\end{itemize}
This structure supports institutional auditing and third-party replication, which is particularly critical in medical settings where pipeline changes can materially alter safety behavior \cite{vollmer2020mlhealth}.

\subsection{Deterministic distillation and traceable supervision}
Distillation pipelines often hide major degrees of freedom in teacher decoding and soft-label generation. MedSwin reduces ambiguity by:
\begin{itemize}
    \item logging the teacher identifier, prompt templates, decoding parameters, and random seeds for hard-label generation \cite{google2025medgemma};
    \item storing top-$k$ token log-probabilities per step for soft KD (renormalized), enabling exact reproduction without re-querying the teacher \cite{hinton2015distilling};
    \item releasing the exact data mixing ratios and ablation configs (SFT-only, KD-only, sequential fine-tuning, training-free merges).
\end{itemize}
This makes MedSwin suitable for controlled comparisons, where differences are attributable to \emph{methods} rather than implicit teacher regeneration.

\subsection{Environment capture and run manifests}
To prevent environment drift, we provide pinned environments (Docker/conda) with exact library versions and hardware notes. Each experiment includes a run manifest documenting: git commit, dataset hashes, GPU model/count, effective batch size, gradient accumulation, training steps, and wall-clock timing. For IR components, manifests include ANN parameters and reranker inference settings, allowing evaluation reproducibility across machines and labs.

\subsection{Reproduction recipe and minimum-cost replication}
A single-command reproduction workflow will rebuild: (i) the SFT checkpoint, (ii) the KD checkpoint (hard + soft), (iii) merged checkpoints using MergeKit YAMLs, and (iv) the evaluation tables under fixed prompts and decoding. This design targets \emph{minimum-cost replication}: independent groups can validate claims without proprietary infrastructure, while preserving enough metadata to support rigorous auditing and extensions.


\section{Appendix}
\label{sec:appendix}

\subsection{Baselines and Ablation Variants}
To isolate the empirical gains of MedSwin, we evaluate a series of controlled ablations alongside representative biomedical baselines across MedQuAD and HealthBench. These configurations are categorized into three tiers:

\begin{enumerate}
    \item \textbf{Reference Tiers:} We include the \emph{base} student (\textit{MedAlpaca-7B}) and the \emph{teacher} model (\textit{MedGemma-27B-IT}) to establish performance bounds for the distillation process.
    \item \textbf{Methodological Ablations:} To decouple the effects of our training objectives, we evaluate \emph{single-objective} variants-specifically \emph{KD-only} (distillation) and \emph{SFT-only} (instruction tuning), and a \emph{sequential} baseline (\emph{SFT-on-KD}) representing standard continued fine-tuning.
    \item \textbf{Merging Heuristics:} We compare our approach against training-free composition methods, including \textsc{Task Arithmetic}, \textsc{Ties}, \textsc{NuSL}, \textsc{Dare-Linear}, and \textsc{Dare-Ties}. To ensure a fair comparison, we maintain a consistent scaling regime across all merge operators to isolate the effect of the merging logic from update magnitude.
\end{enumerate}

For broader external positioning, we benchmark MedSwin against established medical domain models, including \textit{MedGemma-1.5}, \textit{BioMistral-7B}, and \textit{Meditron-7B}, as well as the frontier general-purpose baselines described in the evaluation design.


\subsection{SFT Data Sources, Augmentation, and Quality Control}
\label{app:sft-data-augmentation}

We construct the SFT corpus from complementary public sources to cover the dominant linguistic regimes encountered by clinical assistants: patient-facing symptom narratives, clinician-style explanations and differential reasoning, and evidence-oriented biomedical QA grounded in scientific literature. This mixture reduces distributional bias that arises when training on only one genre.

Supervision is assembled from HealthCareMagic ($\sim$100k) and iCliniq ($\sim$10k)\footnote{HealthCareMagic and iCliniq datasets are publicly available via the ChatDoctor project: \url{https://github.com/Kent0n-Li/ChatDoctor}}\cite{li2023chatdoctor} to capture patient-facing complaint narratives and pragmatic advice-seeking behaviour, and from PubMedQA \cite{pubmedqa} to inject evidence-oriented biomedical reasoning grounded in scientific abstracts. PubMedQA includes expert-labeled, unlabeled, and synthesized splits; the unlabeled split is suitable for teacher-generated hard labels, while the synthesized split supports controlled multi-variant instruction formats.

To enhance coverage and robustness, we apply medically constrained LLM paraphrasing, back-translation, style standardisation, multi-variant question-answer generation, scenario conditioning, and quality assurance. Quality control includes PHI scrubbing, MD5-based deduplication, invalid-generation filtering, and LLM-based medical consistency checks. Augmentation is used to broaden linguistic coverage rather than invent new medical facts. We reject transformations that alter dosage, contraindications, temporal qualifiers, or certainty statements.

\begin{table}[h!]
\centering
\caption{Training corpora and augmentation coverage.}
\label{tab:data}
\begin{tabular}{lccc}
\toprule
\scriptsize
\textbf{Source} & \textbf{Raw Size} & \textbf{Aug.\ Ratio} & \textbf{Notes} \\
\midrule
HealthCareMagic & $\sim$100k & 30-60\% & forum QA \\
HealthCareMagic & $\sim$100k & KD\% & teacher soft/hard labels \\
iCliniq & $\sim$10k & 30-60\% & forum QA \\
PubMedQA (lab.) & 1k & 20-40\% & research QA \\
PubMedQA (unlab.) & 61.2k & KD & teacher hard labels \\
Synthesized & 211.3k & N/A & multi-variant \\
\bottomrule
\end{tabular}
\vspace{-0.5em}
\end{table}


\subsection{Distillation Data Flow}
\label{app:kd-dataflow}

We prioritise unlabeled and synthesized prompts for teacher completion because they provide high diversity of medical intents and question forms without requiring additional human annotation. Teacher hard labels convert these prompts into instruction-following targets, while soft labels preserve calibration and uncertainty structure from the instructor. This is important in medical settings where safe responses often require hedging, uncertainty expression, and evidence qualification.

\begin{figure}[h!]
    \centering
    \includegraphics[width=1\linewidth]{distillation.png}
    \caption{Distillation and adaptation architecture for MedSwin-KD.}
    \label{fig:distillation-pipeline}
\end{figure}


\subsection{Distillation Practicals}
\label{app:kd-practicals}

To transfer both answer content and teacher uncertainty, MedSwin combines hard-label and soft-label distillation. For hard labels, teacher-generated completions are treated as standard instruction targets and mixed with the original SFT corpus. For soft labels, we record the teacher token distribution and train the student to match it.

Let $y_t$ denote the next token at decoding step $t$, $p_T(\cdot|\tau)$ the teacher distribution at temperature $\tau$, and $p_S(\cdot)$ the student distribution. The per-token objective is

\begin{equation}
\mathcal{L}_t
=
\alpha
\big(
-\log p_S(y_t)
\big)
+
(1-\alpha)\tau^2
D_{\mathrm{KL}}
\Big(
p_T(\cdot|\tau)
\,\|\,
p_S(\cdot)
\Big),
\end{equation}

where $\alpha$ controls the balance between hard-label supervision and distribution matching.

To reduce storage cost, only the top-$k$ teacher probabilities are retained and renormalised before training.

\begin{table}[h!]
\centering
\caption{Distillation and adaptation settings.}
\label{tab:settings}
\begin{tabular}{lc}
\toprule
\textbf{Component} & \textbf{Setting} \\
\midrule
Teacher & MedGemma-27B-it \cite{google2025medgemma} \\
Student & MedAlpaca-7B \cite{han2023medalpaca} \\
KD & top-$k$ per-step soft labels; hard labels on unlabeled/synth. \\
Loss & $\alpha$ CE $+$ $(1\!-\!\alpha)\tau^2$ KL, $\tau\!\in\![1,4]$ \\
QLoRA & 4-bit NF4, LoRA rank $r\!\in\!\{8,16\}$, chkpt.\ enabled \\
Regularization & early stopping; cosine/linear warmup \\
\bottomrule
\end{tabular}
\vspace{-0.5em}
\end{table}

Student training uses QLoRA \cite{dettmers2023qlora}: the frozen backbone is loaded in 4-bit NF4 quantisation with low-rank adapters, gradient checkpointing, and AdamW. Early stopping on held-out clinical QA controls overfitting to synthetic or teacher-specific idiosyncrasies, while cosine or linear warmup schedules stabilise joint SFT+KD optimisation.


\subsection{Merge Operators and Technical Rationale}
\label{app:merge-operators}

Unless stated otherwise, a global scaling coefficient is applied to control update magnitude and avoid over-steering the base.

\paragraph{Task Arithmetic}
Task Arithmetic composes deltas additively:
\begin{equation}
\theta_{\text{TA}} = \theta_0 + \lambda_{\text{SFT}}\Delta_{\text{SFT}} + \lambda_{\text{KD}}\Delta_{\text{KD}}.
\end{equation}
It is used as a baseline because it can work when deltas encode complementary behaviours, but it may amplify interference when SFT and KD updates disagree in sign or support \cite{ilharco2022editing}.

\paragraph{SLERP / NuSLERP}
Spherical linear interpolation interpolates between normalised directions:
\begin{equation}
\begin{split}
\Delta(\beta) = & \frac{\sin((1-\beta)\Omega)}{\sin(\Omega)}\widehat{\Delta}_{\text{SFT}}+\frac{\sin(\beta\Omega)} {\sin(\Omega)}\widehat{\Delta}_{\text{KD}}, \\
& \Omega=\arccos(\widehat{\Delta}_{\text{SFT}}\cdot\widehat{\Delta}_{\text{KD}}).
\end{split}
\end{equation}
The merged model is $\theta_{\text{SLERP}}=\theta_0+\gamma\Delta(\beta)$. NuSLERP normalises deltas at a finer granularity, such as per-layer, to reduce layer-wise scale imbalance.

\paragraph{DARE}
DARE sparsifies delta parameters by dropping a fraction $p$ of entries and rescaling the survivors:
\begin{equation}
\widetilde{\Delta}=\mathrm{DARE}(\Delta;p)=\frac{\mathbf{m}\odot \Delta}{1-p}, \quad m_i\sim\mathrm{Bernoulli}(1-p).
\end{equation}
This reduces the probability that conflicting coordinates from SFT and KD are simultaneously active while preserving the expected update $\mathbb{E}[\widetilde{\Delta}]=\Delta$ \cite{yu2023supermario}.

\paragraph{TIES}
TIES targets low-signal updates and sign disagreement by trimming low-magnitude entries, electing a consensus sign per coordinate, and merging only entries whose signs agree with the consensus \cite{yadav2023ties}. This changes the merge from averaging all deltas to composing aligned, high-signal coordinates.

\paragraph{DARE-TIES and DARE-Linear}
DARE-TIES combines stochastic sparsification with sign-aware selective merging. DARE reduces dense overlap, while TIES prevents remaining overlap from cancelling through sign conflict. DARE-Linear applies DARE followed by simple linear fusion, isolating the effect of sparsification without explicit sign-consensus filtering.


\subsection{Biomedical Reranker Training Data and Curation}
\label{app:reranker-data-curation}

We select \emph{open, English} supervision to ensure licensing clarity, reproducibility, and compatibility with local deployment constraints. The reranker must generalise across scientific writing, biomedical information needs, and long heterogeneous passages resembling guideline or EMR-adjacent evidence. We therefore combine three complementary sources.

\begin{enumerate}
    \item \textbf{SciDocs (BEIR/scidocs)} provides relevance supervision derived from citation and qrels-style annotations. Although not strictly clinical, it acts as a scientific ranking prior that improves robustness to literature-grounded retrieval, including PubMed-style abstracts and evidence summaries.
    \item \textbf{BioASQ-generated-queries} provides biomedical query distributions at scale, including abbreviations, entity-heavy phrasing, symptom terms, disease terms, and vocabulary-overlapping hard negatives.
    \item \textbf{MIRIAD 4.4M} contributes broad medical-domain coverage and longer passages, increasing exposure to multi-paragraph evidence where relevance depends on recommendations, contraindications, and conditional statements.
\end{enumerate}

Each training record provides $(q,P,N)$ with one or more positives $P$ and negatives $N$. To mitigate shortcut learning and dataset-specific artefacts, we de-duplicate near-duplicates, preserve section structure during truncation, interleave sources using a fixed mixture schedule, and mine hard negatives from the stage-1 retrieval distribution so that training candidates better match deployment-time candidates.

\subsection{Biomedical Reranker}
\label{subsec:biomedical-reranker}

\textbf{Motivation.}
While domain-specific bi-encoders for biomedical retrieval are increasingly available, open, long-context biomedical rerankers with calibrated, policy-usable scores suitable for auditable EMR/CPG evidence selection remain limited in the literature. Safety-critical clinical RAG requires a cross-encoder that (i) handles multi-paragraph inputs, (ii) exposes a scalar relevance score per pair $(q,d)$ for auditability, and (iii) can be adapted with lightweight PEFT to institution constraints. To close this gap, we adapt the LLM-based cross-encoder \emph{BAAI/bge-reranker-v2-gemma} into a \emph{biomedical reranker} via LoRA finetuning on \emph{open, English} supervision.

\paragraph{Base model and adaptation}
We employ a decoder-only \emph{pointwise} LLM reranker where each query–passage pair $[q;d]$ is scored independently. The model outputs a logit $\ell_\theta(q,d)\in\mathbb{R}$ and an uncalibrated probability $p_\theta(q,d)=\sigma(\ell_\theta(q,d))$, with $\sigma(x)=1/(1+e^{-x})$. We attach LoRA adapters to attention and MLP projections to enable parameter-efficient finetuning under limited VRAM, supporting \textit{query\_max\_len=256} and \textit{passage\_max\_len=1024} tokens (mixed precision; gradient checkpointing).

\paragraph{Training supervision}
We finetune the biomedical reranker using open, English scientific and biomedical retrieval supervision spanning literature-style passages, biomedical question answering, and long guideline-like evidence. The training mixture combines SciDocs, BioASQ-generated queries, and MIRIAD to expose the model to scientific relevance, entity-heavy biomedical queries, and multi-paragraph medical evidence \cite{beirscidocs,bioasqgen,miriad}. Each training record provides $(q,P,N)$ with one or more positives $P$ and negatives $N$. Dataset-specific roles, curation rules, source interleaving, and hard-negative construction are detailed in Appendix~\ref{app:reranker-data-curation}.


\paragraph{Finetuning objective}
We optimise a pointwise binary cross-entropy (BCE) over pairs:
\begin{equation}
\label{eq:bce-reranker}
    \mathcal{L}_{\mathrm{BCE}}
    \;=\;
    -\!\!\sum_{d\in P}\!\log \sigma(\ell_\theta(q,d))
    \;-\!\!\sum_{d\in N}\!\log\big(1-\sigma(\ell_\theta(q,d))\big).
\end{equation}
Triplet-style hardness is induced implicitly by the rolling negative buffer and optional mined negatives from the stage-1 retriever.

\paragraph{Calibration for policy use}
To enable threshold-based EMR auto-inclusion and guideline sufficiency checks, we calibrate logits on a held-out validation set using Platt/temperature scaling:
\begin{equation}
\label{eq:calib}
    \widehat{p}_\theta(q,d)\;=\;\sigma\!\Big(\tfrac{\ell_\theta(q,d)-b}{T}\Big),
\end{equation}
estimating bias $b$ and temperature $T$ by minimising validation NLL. Calibrated probabilities $\widehat{p}_\theta$ are consumed directly by the evidence inclusion thresholds and the evidence-sufficiency gate in Section~\ref{subsec:retrieval-reranking-architecture}, enabling deterministic and  auditable retrieval policies.


The LLM cross-encoder design supports long-context clinical passages, calibrated pairwise relevance scoring, and lightweight site adaptation through LoRA. We use $q$ for the clinical query, $d$ for a candidate passage, $\ell_\theta$ for the reranker logit, and $\widehat{p}_\theta$ for the calibrated probability in Eq.~\ref{eq:calib}. Additional architectural rationale is provided in Appendix~\ref{app:reranker-design-rationale}.


\subsection{Reranker Design Rationale}
\label{app:reranker-design-rationale}

MedSwin uses an LLM-based cross-encoder reranker because clinical evidence often spans multi-paragraph passages containing recommendations, tables, contraindications, subgroup qualifiers, and practice points. Compared with shorter-context BERT-style rerankers, the LLM cross-encoder provides longer effective context, stronger cross-sentence evidence linkage, and compatibility with parameter-efficient adaptation through LoRA. This makes it suitable for auditable medical RAG settings where each query-passage pair must receive a scalar relevance score that can be calibrated and consumed by downstream sufficiency policies.


\subsection{Retrieval Evidence Hierarchy and Source-validity Features}
\label{app:evidence-hierarchy}

For biomedical treatment recommendation, source identity alone is insufficient. Each passage is assigned an evidence-grade vector
\[
\mathbf{e}(d)=
[
e_{\mathrm{CPG}},
e_{\mathrm{SR}},
e_{\mathrm{RCT}},
e_{\mathrm{OBS}},
e_{\mathrm{CASE}},
e_{\mathrm{EMR}},
e_{\mathrm{SAFETY}}
],
\]
where the components represent guideline evidence, systematic reviews, randomised controlled trials, observational studies, case-level evidence, patient-specific EMR evidence, and safety-specific evidence. The EBM score is computed as
\begin{equation}
\begin{aligned}
g_{\mathrm{EBM}}(d) &= \sum_{r} \pi_r e_r(d), \\
\pi_{\mathrm{CPG}}, \pi_{\mathrm{SR}}, \pi_{\mathrm{RCT}} &> \pi_{\mathrm{OBS}} > \pi_{\mathrm{CASE}}.
\end{aligned}
\end{equation}
Additional penalties are applied for outdated recommendations, population mismatch, lack of outcome relevance, weak provenance, or semantic drift. EMR evidence is not treated as higher-level therapeutic evidence; rather, it modulates patient applicability, contraindication detection, medication constraints, and comorbidity alignment.


\subsection{Retrieval Implementation Details}
\label{app:retrieval-implementation}

In implementation, $\mathrm{LCB}_{1-\delta}[\pi_j(\mathcal{M})]$ is estimated by propagating calibration uncertainty from $\widehat{p}_\theta$, facet-alignment uncertainty from $a_{dj}$, and evidence-grade uncertainty from $g_{\mathrm{EBM}}$. A lightweight implementation may use bootstrap resampling over passages and calibration bins, whereas a fully Bayesian implementation may place beta distributions over $a_{dj}$ and $\widehat{p}_\theta$ and compute the posterior lower quantile of $\pi_j(\mathcal{M})$.

To prevent semantic drift and context dilution, the selected evidence bundle is compressed through facet-aware evidence packing. A candidate passage $d$ is protected from pruning if
\begin{equation}
\scriptsize
\label{eq:protected-pruning}
\small
\Delta U(d\mid\mathcal{M})>\epsilon_{\mathrm{pack}}
\lor
R_{\mathrm{safety}}(d)\ge r_{\mathrm{crit}}
\lor
C(d,\mathcal{M})\ge c_{\mathrm{crit}},
\end{equation}
where $R_{\mathrm{safety}}(d)$ measures high-severity safety relevance and $C(d,\mathcal{M})$ measures whether $d$ introduces a clinically meaningful contradiction with the current bundle. Low-relevance, low-quality, or population-mismatched passages are removed when marginal utility is negligible, while high-severity contraindications, subgroup exclusions, and high-grade contradictory recommendations are routed to the contradiction ledger.

Documents are segmented into coherent, section-aware chunks with overlap chosen to minimise topic drift. Each chunk carries $(\mathrm{doc\_id},\mathrm{section},t_{\mathrm{start/end}},\mathrm{source})$ metadata used throughout retrieval, fusion, policy checks, citation construction, and trace export. Typical ANN settings include HNSW $M\!\in\![16,48]$, $\mathrm{efSearch}\!\in\![64,256]$, IVF-PQ $n_{\mathrm{list}}\!\approx\!\sqrt{|\mathcal{D}|}$, and $n_{\mathrm{probe}}\!\in\![8,64]$. Stage~1 retrieval is sublinear per query with ANN, Stage~2 reranking is $O(|\mathcal{C}|)$ pair scorings, and selection is $O(|\mathcal{C}|\cdot|\mathcal{M}|)$ when redundancy penalties are computed against the growing bundle.


\subsection{Multi-Agent Coordination and Audit Artefacts}
\label{app:agent-artifacts}

This appendix provides the implementation details for the Multi-Agent Coordination (MAC) layer described in Section~\ref{subsec:multi-agent-coordination}. The main body presents MAC as a bounded evidence-exploration mechanism. Here, we define the candidate pooling, evidence ledger, reliability weighting, contradiction handling, retrieve-more policy, and audit outputs.

\subsubsection*{Agent-specific evidence exploration}

Let $\{A_i\}_{i=1}^{N}$ denote the set of role-specialised agents. Each agent receives the clinical query $q$, the current facet set $\mathcal{F}(q)$, and the current evidence state. Agent $A_i$ applies a retrieval policy $\pi_i$ and scoring function $S_i(q,d)$ specific to its role. For example, the guideline agent prioritises CPG passages, while the safety agent searches specifically for contraindications and adverse events.

Agent $A_i$ returns a candidate set
\[
\mathcal{C}_i(q)=\{d_{i1},\ldots,d_{iK_i}\},
\]
together with metadata explaining why each passage was retrieved. MAC forms a pooled candidate set:
\begin{equation}
\label{eq:mac-pool}
\mathcal{C}_{\mathrm{MAC}}(q)
=
\bigcup_{i=1}^{N}
\left\{
(d,S_i(q,d),A_i)
\mid
d\in\mathcal{C}_i(q)
\right\}.
\end{equation}

This pooling step is separated from final evidence selection. MAC does not simply take a weighted top-$K$ union, because doing so may suppress rare but safety-critical evidence. Final selection remains controlled by the evidence-sufficiency gate.

\subsubsection*{Structured evidence ledger}

Each agent emits a claim-level evidence ledger rather than an unstructured summary. For agent $A_i$, the ledger is:
\[
\mathcal{L}_i
=
\{
(d,f_j,S_i(q,d),a^{(i)}_{dj},g_{\mathrm{EBM}}(d),
\mathrm{polarity},\mathrm{provenance})
\}.
\]
Here, $f_j\in\mathcal{F}(q)$ is the clinical facet addressed by passage $d$, $a^{(i)}_{dj}$ is the agent's facet-alignment estimate, $g_{\mathrm{EBM}}(d)$ is the evidence-quality score, and $\mathrm{polarity}$ indicates whether the passage supports, contradicts, qualifies, or fails to support a candidate clinical claim.

A treatment-relevant claim is represented as:
\[
c =
(f_j,\mathrm{claim},\mathrm{polarity},d,g_{\mathrm{EBM}},
\widehat{p}_{\theta},\mathrm{provenance}).
\]
This representation links each claim to a clinical facet, source passage, evidence grade, calibrated relevance score, and provenance record.

\subsubsection*{Reliability-weighted aggregation}

Each agent is assigned a reliability weight estimated from validation-time performance. Let $r_i$ denote the probability that agent $A_i$ contributes clinically valid, non-redundant, and evidence-policy-compliant information for the current query class. We model:
\[
r_i\sim\mathrm{Beta}(a_i,b_i),
\]
where $(a_i,b_i)$ are updated from validation signals such as retrieval precision, safety recall, calibration error, contradiction-resolution accuracy, and policy-compliance rate.

To avoid over-weighting agents with uncertain estimates, MAC uses a conservative lower confidence bound:
\[
\widetilde{r}_i
=
\mathrm{LCB}_{1-\delta}(r_i),
\qquad
\omega_i
=
\frac{\widetilde{r}_i}
{\sum_{k=1}^{N}\widetilde{r}_k}.
\]
Thus, an agent's influence depends on demonstrated reliability rather than its assigned role alone.

\subsubsection*{Reliability-weighted bundle selection}

The final MAC evidence bundle is selected from the pooled candidate set by maximising a reliability-weighted clinical utility under a token budget:
\begin{equation}
\label{eq:mac-selection}
\mathcal{M}_{\mathrm{MAC}}^{\star}
=
\arg\max_{\mathcal{M}\subseteq\mathcal{C}_{\mathrm{MAC}}(q)}
U_{\mathrm{MAC}}(q,\mathcal{M})
\quad
\text{s.t.}
\quad
\sum_{d\in\mathcal{M}}\tau(d)\le B.
\end{equation}

The utility is defined as:
\begin{equation}
\label{eq:mac-utility}
\scriptsize
\begin{aligned}
U_{\mathrm{MAC}}(q,\mathcal{M}) 
&= \sum_{d\in\mathcal{M}} \left[ \sum_{i=1}^{N}\omega_i S_i(q,d) + \lambda_{\mathrm{agree}}A(d) \right. \\
&\qquad \left. + \,\lambda_{\mathrm{EBM}}g_{\mathrm{EBM}}(d) + \lambda_{\mathrm{safety}}R_{\mathrm{safety}}(d) \right] \\
&\quad - \lambda_{\mathrm{red}}R_{\mathrm{red}}(\mathcal{M}) - \lambda_{\mathrm{conf}}R_{\mathrm{conf}}(\mathcal{M}) - \lambda_{\mathrm{noise}}R_{\mathrm{noise}}(\mathcal{M}).
\end{aligned}
\end{equation}

Here, $A(d)$ rewards independent corroboration by reliable agents, $R_{\mathrm{safety}}(d)$ rewards evidence about contraindications and high-severity risks, $R_{\mathrm{red}}$ penalises redundant evidence, $R_{\mathrm{conf}}$ penalises unresolved contradiction among high-quality sources, and $R_{\mathrm{noise}}$ penalises population mismatch, obsolete recommendations, weak provenance, or semantic drift.

Final acceptance remains governed by the evidence-sufficiency gate. A high MAC utility score does not automatically permit answer generation if a critical facet remains missing or a high-severity contradiction remains unresolved.

\subsubsection*{Contradiction adjudication}

Contradictions are not averaged away during aggregation. If two high-quality passages make incompatible claims about the same clinical facet, MAC records a contradiction pair:
\[
(d_a,d_b,f_j,\mathrm{severity},g_{\mathrm{EBM}}(d_a),g_{\mathrm{EBM}}(d_b)).
\]

Low-severity contradictions may be resolved using evidence hierarchy, recency, source reliability, or population fit. High-severity contradictions, especially those involving contraindications, adverse events, medication restrictions, or guideline-discordant recommendations, are preserved in the contradiction ledger and routed to the contradiction critic. If the conflict remains unresolved, the final answer must state the uncertainty rather than presenting a single unsupported recommendation.

\subsubsection*{Targeted retrieve-more actions}

When the sufficiency gate is not satisfied, MAC issues targeted \textit{retrieve-more} actions. The missing facet determines the next retrieval route:

\begin{itemize}
    \item missing or weak guideline evidence triggers the guideline agent;
    \item missing patient applicability triggers the EMR agent;
    \item missing contraindication or adverse-risk evidence triggers the safety agent;
    \item weak evidence grading triggers the evidence-quality agent;
    \item unresolved source conflict triggers the contradiction critic.
\end{itemize}

This prevents the system from repeatedly expanding the context window without a clear purpose. Additional retrieval is justified only when it is expected to improve coverage of a missing or uncertain clinical facet.

\subsubsection*{Audit artefacts}

For each query, MAC exports the following audit artefacts:

\begin{enumerate}
    \item \textbf{Retrieval trace:} records which agent retrieved each passage, from which source, and for which facet.
    \item \textbf{Reranking trace:} records calibrated relevance scores and selected/rejected candidate passages.
    \item \textbf{Facet-coverage matrix:} records which passages support each clinical facet and whether each facet is covered, partial, missing, or contradicted.
    \item \textbf{Contradiction ledger:} records high-severity conflicts, source pairs, affected facets, severity, and resolution status.
    \item \textbf{Sufficiency decision:} records whether the system answered, retrieved more evidence, or returned an insufficient-evidence response.
    \item \textbf{Final-answer provenance:} links each treatment-relevant statement in the final response to accepted evidence passages.
\end{enumerate}

These artefacts make the stopping decision inspectable. A reviewer can determine whether MedSwin answered because the evidence bundle satisfied the sufficiency gate, retrieved more because a critical facet was missing, or limited the answer because the available evidence did not support a full recommendation.


\subsection{Benchmark data selection and preparation}
We separate model-level and system-level evaluation because they test different failure modes. MedQuAD and HealthBench are retained as model-level benchmarks because they measure whether a model can produce semantically aligned medical answers under a fixed prompting and scoring pipeline. These benchmarks are appropriate for comparing the MedSwin student, teacher, SFT, KD, and merge variants against external biomedical and frontier baselines. However, answer-only benchmarks cannot determine whether a clinical RAG system retrieved the correct evidence, preserved safety-critical passages, applied sufficiency gates, or exposed an auditable decision trace.

For the whole-system audit, we use TREC Clinical Decision Support 2016 as the primary benchmark because it more closely matches MedSwin's intended operating mode: patient-context clinical decision support with retrieval from a biomedical evidence corpus and judged relevant documents. Each case provides a clinical topic derived from patient context, and the released relevance judgments allow evidence-document recall and citation precision to be measured without relying solely on generated text similarity. This makes TREC-CDS a stronger test of the integrated retrieval, reranking, evidence-selection, sufficiency, and trace-generation pipeline.

In the prepared case file, the clinical query is kept separate from the patient note. The patient note is ingested as scoped EMR evidence, while biomedical articles are ingested as literature evidence. This design prevents the benchmark from becoming a prompt-only summarisation task and forces the runtime to use its normal retrieval and evidence-selection contract. It also allows the audit to distinguish three clinically important behaviours: retrieving patient-specific context, retrieving external biomedical evidence, and deciding whether the resulting bundle is sufficient for a bounded clinician decision-support response.


\subsection{Metric Definitions}
\label{app:metric-definitions}

For model-level evaluation, ROUGE-L F1 measures sequence-level overlap using longest common subsequence matching \cite{lin2004rouge}. BERTScore-F1 measures semantic similarity under contextual embeddings and is more tolerant to paraphrase than n-gram overlap \cite{bertscore}. Token-level F1 is used as a stricter lexical correctness proxy, while unigram and bigram precision are used as over-generation proxies. We treat n-gram precision as a proxy rather than a definitive hallucination metric because faithfulness and hallucination measurement remain metric-dependent \cite{maynez-etal-2020-faithfulness,ji2023hallucination,kulkarni-etal-2025-evaluating,mickus-etal-2024-semeval}.

For the system audit, \textbf{facet recall} measures the proportion of required clinical facets supported by the selected evidence bundle. \textbf{Critical facet recall} restricts this calculation to safety-critical or decision-critical facets, such as contraindications, treatment-relevant exclusions, adverse-risk evidence, and patient applicability. \textbf{Evidence-document recall} compares selected or cited document identifiers against benchmark relevance judgments. \textbf{Citation precision} measures the proportion of cited documents judged relevant.

\textbf{Trace completeness} measures whether the returned audit artifact contains the required review fields, including policy decision, selected evidence, citations, facet coverage, evidence ledger, contradiction ledger, and degraded-mode state. \textbf{Sufficiency decision score} evaluates whether the generation gate is consistent with the evidence state. \textbf{Groundedness proxy} estimates whether answer content is supported by the selected evidence bundle. \textbf{Unsupported penalty} penalises clinical claims not supported by selected evidence, while \textbf{unsafe omission penalty} penalises missing critical facets.


\subsection{Composite Benchmark Indices}
\label{app:composite-benchmark-indices}

The raw benchmark metrics remain the primary results. The composite indices are used only to make the figure easier to read when comparing more than twenty models across five metrics. This follows common multi-metric benchmark practice, where aggregate views are used for compact comparison while metric-level results are still reported separately \cite{wang2018glue,wang2019superglue,liang2023holistic,muennighoff2023mteb}. We also follow standard composite-indicator practice by making the normalization and weighting explicit \cite{nardo2008handbook}.

For each dataset, every metric is min-max normalized independently:
\[
\tilde{m}_{i}
=
\frac{m_i - \min_j m_j}{\max_j m_j - \min_j m_j}.
\]
The \textbf{Answer Quality Index} is:
\begin{equation}
\scriptsize
\begin{aligned}
\mathrm{AQI}_i = 100 \times \Big( &0.60 \widetilde{\mathrm{BERTScore}}_i + 0.25 \widetilde{\mathrm{TokenF1}}_i + 0.15 \widetilde{\mathrm{ROUGE\text{-}L}}_i \Big).
\end{aligned}
\end{equation}
BERTScore receives the largest weight because it better captures semantic similarity under paraphrase than exact lexical overlap \cite{bertscore}.

The \textbf{Reference Precision Index} is:
\begin{equation}
\scriptsize
\begin{aligned}
\mathrm{RPI}_i = 100 \times \Big( &0.60 \widetilde{\mathrm{BigramPrecision}}_i + 0.40 \widetilde{\mathrm{UnigramPrecision}}_i \Big).
\end{aligned}
\end{equation}

We use the term reference precision rather than hallucination score because n-gram precision cannot directly establish factuality or clinical grounding \cite{maynez-etal-2020-faithfulness,ji2023hallucination}.

The ranked panel reports:
\[
\mathrm{BCI}_i
=
0.70 \mathrm{AQI}_i
+
0.30 \mathrm{RPI}_i.
\]
This Benchmark Composite Index is a visualization aid only. Clinical reliability is evaluated separately using evidence sufficiency, critical-facet recall, citation precision, and trace completeness.


\subsection{Retrieval-style reranker benchmark}
For MedSwin-Rerank, each MedQuAD example is converted into a ranking problem where the question is the query and the reference answer is the positive passage. Candidate lists are built with one positive and a fixed number of negatives, either sampled randomly or retrieved first and reranked. The benchmark then reports MRR@10, nDCG@10, Recall@10, MAP@10, and p95 latency.


\subsection{TREC CDS case preparation}
The case-preparation script exports a compact JSONL file from the TREC-CDS loader. Each record contains a case identifier, query type, clinical information need, patient-context note, gold relevant document identifiers, and seeded audit facets. Gold evidence is derived from the TREC qrels, with higher relevance grades prioritised before lower grades. When a topic contains more judged documents than the configured cap, document ID is used as a stable tiebreaker. The benchmark therefore avoids random negative sampling in the system-level audit and remains exactly reproducible across runs.

Facet labels are seeded from the query type. Diagnosis cases require diagnostic evidence and patient applicability; test cases require test indication and patient applicability; treatment cases require treatment recommendation evidence, safety or adverse-risk evidence, and patient applicability. This facet design intentionally reflects MedSwin's sufficiency policy. The goal is not to ask whether the model can produce a plausible paragraph, but whether the evidence bundle contains the minimum information needed to support a safe clinical decision-support response.

The companion ingestion script loads PMC evidence into a fixed benchmark organisation namespace and ingests each patient note as EMR evidence scoped to its case-specific patient identifier. Literature evidence remains organisation-scoped rather than patient-scoped, preventing the patient filter from excluding external biomedical documents. After ingestion, embeddings are refreshed and the index is rebuilt so dense retrieval, optional lexical retrieval, reranking, and evidence selection operate over the same evidence pool.


\subsection{Detailed Worked Example Trace}
\label{app:worked-example-trace}

This appendix expands the medication-safety example from Section~\ref{subsec:worked-example}. The query is:

\begin{quote}
\emph{Can this patient continue metformin after the latest renal-function result?}
\end{quote}

The example is used only to illustrate the MedSwin workflow. It does not provide clinical advice.

\paragraph{Step 1: Facet identification.}
MedSwin decomposes the query into six clinical evidence facets:
\begin{enumerate}
    \item current metformin use;
    \item latest renal-function status;
    \item renal-function trend or acute-change context;
    \item relevant CPG threshold or caution statement;
    \item medication-safety warning;
    \item unresolved contradiction among sources.
\end{enumerate}

\paragraph{Step 2: Candidate retrieval.}
The EMR agent retrieves medication lists, recent laboratory results, and historical renal-function records. The guideline agent retrieves CPG passages related to metformin use in renal impairment. The safety agent retrieves drug-safety passages describing relevant adverse-risk concerns.

\paragraph{Step 3: Reranking.}
The biomedical reranker scores each candidate passage for query relevance. Passages directly affecting the medication-safety decision, such as the latest renal result and the guideline renal-use threshold, are ranked above general diabetes-management background.

\begin{table}[h!]
\centering
\caption{Illustrative reranking output for the worked example.}
\label{tab:appendix-worked-example-rerank}
\scriptsize
\begin{tabularx}{\linewidth}{l X l c}
\toprule
\textbf{ID} & \textbf{Evidence role} & \textbf{Source} & \textbf{Score} \\
\midrule
EMR-1 & Confirms active metformin use & EMR medication list & 0.88 \\
EMR-2 & Provides latest renal-function result & EMR lab result & 0.94 \\
CPG-1 & Gives renal-use threshold or caution statement & CPG & 0.91 \\
SAF-1 & Describes medication-safety risk & Drug-safety source & 0.80 \\
EMR-3 & Historical renal-function trend unavailable or incomplete & EMR history & - \\
\bottomrule
\end{tabularx}
\end{table}

\paragraph{Step 4: Sufficiency decision.}
The evidence bundle covers current medication, latest renal function, guideline threshold, and safety warning. However, the renal trend or acute-change context remains missing. Since this facet affects whether the latest result should be interpreted as stable chronic impairment or acute deterioration, the system marks the bundle as partially insufficient.

\paragraph{Step 5: Policy action.}
MedSwin issues one of two allowed actions:
\begin{enumerate}
    \item retrieve additional laboratory history or clinical notes related to renal-function change; or
    \item return a bounded insufficient-evidence response that states the missing trend evidence and avoids a confident medication recommendation.
\end{enumerate}

\paragraph{Step 6: Final bounded response.}
If additional retrieval fails, the response should be constrained:

\begin{quote}
The available evidence confirms current metformin use, the latest renal-function result, and a relevant guideline caution. However, the renal-function trend or acute-change context is missing. Therefore, the system cannot determine from the available evidence alone whether the result reflects stable chronic impairment or acute deterioration. The medication decision should be reviewed against the guideline threshold once the missing renal trend and clinical context are confirmed.
\end{quote}

This trace demonstrates the intended behaviour of MedSwin's evidence-sufficiency gate: the system may provide a limited, evidence-grounded response, but it should not present a full recommendation when a critical facet is missing.


\subsection{System audit workflow}
Each system run follows the same sequence. First, the harness performs a health check against the MedSwin runtime. Second, it optionally ingests the case-specific patient note as EMR evidence. Third, it calls the production \texttt{/chat} endpoint using the same request contract used by deployment clients. Fourth, it retrieves the generated trace from \texttt{/traces/\{trace\_id\}}. Finally, it writes an audit JSON record containing the answer, selected passages, cited passages, policy decision, facet coverage, contradiction ledger, evidence ledger, trace summary, degraded-mode flags, and scoring outputs.

This workflow is intentionally endpoint-level rather than component-level. The benchmark does not call the retriever, reranker, or policy module directly. As a result, failures in document scoping, retrieval visibility, reranker calibration, evidence packing, policy thresholds, trace generation, or citation construction appear in the final audit. This design is important for evaluating multi-agent clinical IR systems because isolated component scores can hide integration failures. A system that has a strong reranker but fails to pass selected evidence into the sufficiency policy should score poorly, as should a fluent generator that produces claims without traceable evidence.

The run-level summary reports the mean over cases for each metric, plus policy pass rate, degraded rate, and error rate. Policy pass rate is interpreted jointly with critical-facet recall and unsafe omission penalty: a high pass rate is desirable only when the system is retrieving sufficient evidence, whereas a high pass rate with low critical-facet recall would indicate unsafe over-generation. Conversely, a low pass rate with high trace completeness may indicate safe refusal behaviour but inadequate retrieval. This paired interpretation prevents the benchmark from rewarding either reckless answering or excessive abstention.


\subsection{System audit metrics}
The audit reports mean facet recall, critical facet recall, evidence-document recall, citation precision, trace completeness, sufficiency decision score, groundedness proxy, clinical quality proxy, unsupported penalty, unsafe omission penalty, and the composite MSAS score. Each metric is computed from the production response and the benchmark gold record.

\textbf{Facet recall} measures the proportion of required clinical facets supported by the selected evidence bundle. \textbf{Critical facet recall} is the same calculation restricted to facets marked safety-critical or decision-critical, such as treatment recommendation evidence, contraindications, adverse-risk evidence, and patient applicability. This metric receives the largest weight in MSAS because clinical decision support is especially sensitive to missing pivotal information. A system that retrieves broad background information but misses the contraindication or patient-specific exclusion should not be considered reliable.

\textbf{Evidence-document recall} compares selected or cited document identifiers against the TREC relevance judgments. It measures whether the system retrieves externally judged biomedical evidence rather than only the patient note or generic background passages. \textbf{Citation precision} measures the proportion of cited documents that are relevant according to the benchmark judgments. These two metrics separate evidence recall from citation selectivity: high evidence-document recall with low citation precision indicates noisy evidence selection, whereas high citation precision with low evidence-document recall indicates overly narrow citation behaviour.

\textbf{Trace completeness} measures whether the returned audit artifact contains the fields required for clinical review, including policy decision, selected evidence, citations, facet coverage, evidence ledger, contradiction ledger, and degraded-mode state. \textbf{Sufficiency decision score} evaluates whether the system's generation gate is consistent with the evidence state. It rewards full answers when critical evidence is present and bounded insufficient-evidence responses when critical evidence is missing. This is essential for MedSwin because safe abstention is a valid clinical decision-support behaviour.

\textbf{Groundedness proxy} estimates whether answer content is supported by the selected evidence bundle. \textbf{Clinical quality proxy} captures whether the answer follows the expected decision-support format, including uncertainty, safety notes, and non-autonomous diagnosis boundaries. \textbf{Unsupported penalty} penalises clinical claims not supported by selected evidence, while \textbf{unsafe omission penalty} penalises missing critical facets. These penalties prevent a fluent but weakly grounded answer from scoring highly.

The composite MSAS is therefore not a generic text-generation score. It is a system contract score: it rewards the pipeline for retrieving relevant evidence, selecting critical facets under budget, making an appropriate sufficiency decision, citing evidence precisely, and exposing a trace. This makes it suitable for evaluating multi-agent clinical IR systems, where the central risk is not only incorrect generation but also failed evidence acquisition, missing provenance, or unsafe policy pass-through.

\subsection{Detailed System-Evaluation Results}
\label{app:detailed-system-results}

Table~\ref{tab:trec_cds_system_audit} shows that MedSwin's reliability is not solely a property of its standalone generator or reranker, but emerges from the interaction between retrieval, reranking, facet sufficiency, and evidence-gated generation. The MedSwin backend obtains the highest policy pass rate (0.85), critical-facet recall (0.82), evidence-document recall (0.53), trace completeness (0.81), and tied-best MSAS (0.84), with zero runtime audit errors.

The OpenAI backend achieves the highest overall facet recall (0.92) and citation precision (0.88), suggesting stronger broad coverage and more selective citation behaviour. However, its lower critical-facet recall (0.78) and evidence-document recall (0.45) indicate that high citation precision alone is insufficient when safety-critical or patient-applicability evidence is missed. The GoogleAI backend performs competitively on evidence-document recall (0.50), but its lower critical-facet recall (0.72) and MSAS (0.75) suggest weaker alignment with the sufficiency policy. In contrast, the MedSwin backend trades some citation precision for stronger retrieval of critical evidence and more complete trace generation.


\subsection{Validity of the system audit for multi-agent clinical IR}
The TREC-CDS system audit is appropriate for MedSwin because it evaluates the full clinical IR loop rather than a single generative endpoint. Multi-agent RAG systems can fail in ways that answer-only benchmarks cannot observe: one agent may retrieve irrelevant evidence, another may summarise patient context incorrectly, the reranker may suppress safety-critical minority evidence, the sufficiency policy may pass an incomplete bundle, or the final answer may cite evidence that does not support its claims. The audit exposes these failures by recording selected evidence, cited evidence, policy decisions, facet coverage, and trace completeness for every case.

The benchmark also avoids rewarding ungrounded fluency. A model that generates a clinically plausible answer without retrieving the relevant TREC-judged documents receives low evidence-document recall and citation precision. A model that retrieves relevant documents but fails to satisfy critical facets receives a lower critical-facet score. A model that cannot produce a reviewable trace receives a lower trace-completeness score. Thus, the evaluation aligns with the practical requirements of clinical decision support: the system must find evidence, select it under constraints, explain why it is sufficient or insufficient, and ground its response.

The benchmark has limitations. TREC relevance judgments are document-level rather than claim-level clinical adjudications, and they are not a substitute for prospective clinician evaluation. The seeded facets are derived from query type and qrels rather than from exhaustive expert annotation of every clinical statement. Consequently, MSAS should be interpreted as an end-to-end evidence and auditability score, not as proof of bedside clinical efficacy. Nevertheless, it is a stronger test for MedSwin's intended contribution than answer-only QA because it evaluates integration reliability, provenance, sufficiency behaviour, and safety-sensitive evidence retrieval under a reproducible patient-context retrieval setting.


\subsection{Sequential SFT-on-KD Analysis}
\label{app:sequential-sft-kd}

Sequential SFT-on-KD degrades semantic alignment on both benchmarks. On MedQuAD, BERTScore drops from $0.8441$ for KD-only to $0.8341/0.8297$ for the two KD$\rightarrow$SFT variants; on HealthBench, it similarly falls from $0.8314$ to $0.8186/0.8088$. The lexical metrics decrease less dramatically than the semantic ones, suggesting that continued instruction tuning can preserve parts of answer form while still overwriting teacher-like distributional behaviour learned during KD. This is compatible with continual fine-tuning dynamics, where later updates induce catastrophic forgetting of previously acquired semantics and calibration \cite{luo2023catastrophic}. In contrast, merging avoids additional gradients and therefore avoids the primary mechanism of forgetting, while still permitting controlled incorporation of SFT-like response discipline.


\section{Acknowledgement}
\bibliographystyle{IEEEtran}
\bibliography{references}
\end{document}
